\bookpart{测度空间}{Measure Spaces}
\label{part:measure-spaces}

\BookPartGuide{%
前面的微积分用实数长度和欧氏体积累计函数值。本部分把底集换成一般对象，先规定哪些集合可以测量、如何在其上赋权，再比较同一空间上的两种赋权，以及不同空间之间的质量传递、联合记录与跨空间成本，最后研究测度序列的各种收敛方式与极限子列。
第~\ref{chap:measure-spaces}章“测度空间”建立一般测度、Borel结构与支持，回用已有的积分构造和乘积积分，给出相对密度、测度分解、全变差、Hellinger、KL与Wasserstein等比较工具，并引入可测映射的推前、测度核、联合测度分解和运输计划。
概率论中的条件期望、独立性、Bayes法则与随机过程扩张将在后文专门讨论。}

\chapter{测度空间}
\label{chap:measure-spaces}

在前面的微积分里，我们累计一个量时，底集总是实轴或欧氏空间：长度、面积、体积已经替我们规定好“一块区域有多大”，积分就是在这套规定上累计函数。
但在机器学习中，需要累计的对象往往不是数：它们可能是词、类别、图、数据样本，或者另一个空间里的点。
这时长度和体积帮不上忙，我们必须自己回答两件事——哪些集合允许被测量，以及测量的结果是多少。
回答这两个问题的工具，就是\emph{可测结构}和\emph{测度}。

先看一个具体场景。
给定一批数据上的两种分布，我们会问：它们是不是一回事？如果不是，差在什么地方？
有时两种分布描述的是同一组对象，只是各部分占的比重不同：比如同一批人按两种口径统计收入。
有时它们根本描述不同空间的对象：比如一张图片和它的类别标签，需要一条规则说明图片上的“质量”怎样传到标签上。
这两类问题构成本章的两条主线。

第一条主线固定同一个空间，只比较上面的两种赋权。
我们要判断一种赋权认为是零的集合，另一种是否也认为是零；能不能找到一个密度函数，用一种赋权把另一种表示出来；以及用一个数概括两者的整体差异时，不同指标各在衡量什么。
第二条主线允许空间不同。
给定一个映射时，输入空间的质量怎样“推”到输出空间；当一个输入不只对应一个确定输出，而是按某种概率分散到若干输出时，这种规则怎样描述；如果还想知道输入与输出是怎样配对的，又要付出什么跨空间的移动代价。

\ref{sec:measure-general-structure}节先打基础，把实轴上已经熟悉的可测集合、测度、积分和乘积推广到一般空间。
\ref{sec:measure-same-space-comparison}节回到第一条主线，比较同一空间上的测度。
\ref{sec:measure-cross-space}节处理第二条主线，研究映射、测度核、联合测度和跨空间运输。
\ref{sec:measure-convergence}节再让测度随序号或时间变化，讨论一列测度怎样才算趋近、在什么条件下一定能找到收敛的子列。

这些内容不是几种互不相干的“测度种类”，而是同一套工具在回答不同问题。
确定映射其实是一种特殊的测度核：它给每个输入指派的输出测度全部集中在一个点上。
全变差、Wasserstein等差异指标也有联系：有的只看同一事件上的质量差，有的还关心两点在空间中隔得多远，我们会在相应位置把关系讲清。
概率测度在本章只是总质量恰好为$1$的测度；条件期望、Bayes公式、随机过程等需要额外概率结构的内容都留给后文，这里只建立可以反复调用的通用对象和运算。

\section{一般测度的结构与积分}
\label{sec:measure-general-structure}

在实轴和欧氏空间上，“一块区域有多大”是由长度、面积、体积直接给出的。
换到一般对象的集合后，这种天然的大小不存在了，我们要自己把测量的规则补回来。
本节做四件准备工作：先说清哪些集合允许被测量、再说明怎样给这些集合配质量；然后引入开集、闭集等拓扑概念在测量中的位置；接着把熟悉的积分搬到一般测度上；最后处理两个空间的联合结构。
这些内容在实轴上都已经出现过，这里不重复证明，重点是看清“换了底集以后，哪些结论原封不动，哪些条件必须重新声明”。

\subsection{可测结构与基本赋权}

拿到一个由任意对象组成的集合后，测量要分两步走，而且两步互不相干。
第一步决定\emph{哪些集合看得见、可以测量}，这是可测结构；第二步决定\emph{看得见的集合各自有多大}，这是测度。
同一套可见集合上可以配完全不同的质量，同一个质量的想法在更粗或更细的可见范围下也要分别讨论。
本小节先建立这两层，再给出点质量、测度混合和零测例外的说法。

\subsubsection{可测范围与子$\sigma$代数}

\paragraph{可测空间与生成结构}

设$X$是我们关心的全部对象，例如若干类别、若干样本或若干状态。
要测量它的子集，先要圈定一个“允许测量的集合族”。
为什么不能直接宣布所有子集都可测？
在实轴上已经看到，一旦允许不可数次地分割、覆盖和取极限，就会出现无法自洽赋长度的集合；因此允许范围必须满足一定封闭性，但又不能大到没有限制。

我们沿用$\sigma$代数的要求：集合族$\mathcal A\subseteq\mathcal P(X)$包含全集$X$，并且对相对补集和可数并封闭。
取补和可数并还自动带来另外几种封闭：空集属于$\mathcal A$，可数交
\[
 \bigcap_nA_n=X\setminus\bigcup_n(X\setminus A_n)
\]
属于$\mathcal A$，两个可测集的差$A\setminus B=A\cap(X\setminus B)$也属于$\mathcal A$。
换句话说，凡是能用可测集合作可数多次集合运算得到的集合，仍然允许测量。
配备了这样一个集合族以后，二元组$(X,\mathcal A)$称为可测空间，$\mathcal A$中的集合称为可测集合。

实际问题里，我们往往只先点名少数几个必须能测量的集合，例如区间、矩形或某些类别，其余集合由它们“生成”。
对一族集合$\mathcal C$，包含它的$\sigma$代数可能有很多，其中最小的一个记为$\sigma(\mathcal C)$。
最小性可以这样理解：幂集$\mathcal P(X)$总是一个包含$\mathcal C$的$\sigma$代数，所有这种$\sigma$代数的交仍然包含$X$且对补集、可数并封闭，所以交出来的最小范围确实存在。
这种“只规定少数起点，再封闭出全部可测集”的做法非常省事：后面验证一个映射是否可测、两项测度是否相同，都只需在起点集合上检查，再由封闭性推广到整个$\sigma$代数。

用一个有限例子体会“可见”的含义。
在$X=\{1,2,3,4\}$上，幂集$\mathcal P(X)$允许区分每一个点；但由划分$\bigl\{\{1,2\},\{3,4\}\bigr\}$生成的$\sigma$代数只有
\[
 \bigl\{\varnothing,\{1,2\},\{3,4\},X\bigr\}.
\]
在这套范围里，“第$1$个点”单独是看不见的，能测量的只有前两个点、后两个点这两个分块。
可见集合的多少，决定了我们能在多大精度上讨论问题。

\paragraph{子$\sigma$代数与同一底集上的粗化}

有时底集不变，但我们主动“眯起眼睛”，只看其中一部分集合。
若$\mathcal G\subseteq\mathcal A$本身也是一个$\sigma$代数，就称它为\term{子σ代数}。
$(X,\mathcal G)$中的对象一个不少，仍是原来的$X$，只是可见的集合变少了。
这种操作叫粗化，它和“把底集缩小成子集$Y\subseteq X$”不同：缩小底集是扔掉一批对象，粗化只是降低对同一批对象的分辨能力。

仍以四点集合为例。
在$\mathcal A=\mathcal P(X)$下四个单点都能单独测量；换成
\[
 \mathcal G=\{\varnothing,\{1,2\},\{3,4\},X\},
\]
就只能称“前两点合起来多重”“后两点合起来多重”，再也无法分开第$1$点和第$2$点。
若原来在$\mathcal A$上有测度$\mu$，只在$\mathcal G$上读取它，记作$\mu|_{\mathcal G}$，这只是缩小赋权规则的定义域，并没有删除任何点。
后面讨论\term{条件积分}时，会专门研究“在只能看到$\mathcal G$的情况下，怎样保留各分块上的累计信息”；这里先把粗化这种结构本身建立起来。

\subsubsection{测度、混合与零测例外}

\paragraph{一般测度及有限性条件}

圈定可见集合以后，就要给每个集合配一个非负的“大小”。
设$\mu$把每个可测集映到$[0,+\infty]$。
要求它满足两条：空集的大小为$0$；把互不相交的可数个集合合并时，总大小等于各部分大小之和，即
\[
 \mu\!\left(\bigcup_{n=1}^{\infty}A_n\right)=\sum_{n=1}^{\infty}\mu(A_n).
\]
这样的$\mu$称为测度，$(X,\mathcal A,\mu)$称为测度空间。
第二条（可数可加）是测量的核心：它保证“分成几块分别称，再合起来称”，结果一致。
底集换成一般对象后，这个直觉与长度完全相同；由它还能立刻得到单调性$A\subseteq B\Rightarrow\mu(A)\leq\mu(B)$和可数次可加，证明与实轴上一样，不再重复。

后面的定理对“总质量有多大”有不同要求，这里先把三个常用名词一次分清，免得以后混淆。
若$\mu(X)<\infty$，称$\mu$为\term{有限测度}；进一步若$\mu(X)=1$，称$\mu$为\term{概率测度}，此时每个集合的质量都在$[0,1]$内。
还有一个更弱、却非常常用的条件叫$\sigma$有限，满足它的测度即\term{σ有限测度}：可以用可数个可测集$E_k$盖住$X$，使每一块的质量都有限。
三者关系是：概率测度一定有限，有限测度一定$\sigma$有限（取$E_1=X$），但反过来不成立。
自然数集上的计数测度给每个点质量$1$、全空间质量无穷，却能被单点集$\{1\},\ldots,\{k\}$有限地盖住；实轴上的勒贝格测度总质量无穷，但区间$[-k,k]$每一块都有限——它们都是$\sigma$有限而不有限。

集合列的连续性沿用前面的结论，但要留意有限性出现在哪里。
集合逐渐扩大时，$A_n\uparrow A$就有$\mu(A_n)\uparrow\mu(A)$，不要求任何有限性；集合逐渐缩小时，$A_n\downarrow A$要得到$\mu(A_n)\downarrow\mu(A)$，必须先有$\mu(A_1)<\infty$。
否则$(n,\infty)\downarrow\varnothing$，而每一项的勒贝格长度都是无穷，无法相减得到零。
请记住：$\sigma$有限不是所有运算的万能前提，密度表示、乘积构造和有界函数积分各自需要的条件，会在用到时分别说明。

\paragraph{加权点测度与\term{可数混合}}

最容易想象的测度，是把质量直接放在若干个点上。
固定一个点$x$（并假设$\{x\}$可测），定义
\[
 \delta_x(A)=\mathbf 1_A(x)=
 \begin{cases}1,&x\in A,\\0,&x\notin A,\end{cases}
\]
称为集中在$x$的点质量：集合包含$x$就得质量$1$，否则得$0$。
它对应“全部重量压在一个点上”的极端赋权。

再取一列点$x_i$和非负权重$w_i\geq0$，把若干点质量加起来：
\[
 \mu=\sum_i w_i\delta_{x_i},\qquad
 \mu(A)=\sum_i w_i\,\mathbf 1_A(x_i).
\]
因为每个点属于哪些不交集合是唯一确定的，两个非负求和可以交换次序，所以$\mu$仍然可数可加；同一位置上的权重自动相加。
权重总和有限时$\mu$是有限测度，总和为$1$时就是概率测度。
这里有个细节：定义本身不要求每个单点都可测，但只要具体计算中用到$\delta_x(\{x\})=1$，就必须先确认$\{x\}\in\mathcal A$。

同样的“加权相加”还可以用在测度上。
设$\mu_i$是同一空间上的非负测度，取非负系数$a_i$，令
\[
 \mu(A)=\sum_i a_i\mu_i(A),
\]
非负级数换序保证$\mu$仍是测度。
若各$\mu_i$都是概率测度且系数和为$1$，得到的就是可数混合：以比例$a_i$在第$i$种随机规律之间随机挑选一种。
一个直观例子是
\[
 \mu=\tfrac12\,\lambda|_{[0,1]}+\tfrac12\delta_0:
\]
一半重量像水一样均匀铺在$[0,1]$上按长度分配，另一半重量全部压在点$0$上。
这样的混合提醒我们：一个测度完全可以同时拥有“连续铺开”的部分和“集中成点”的部分。

\paragraph{零测例外与完备化}

测量时，质量为零的集合允许我们忽略其中点上的例外。
可测集$N$若满足$\mu(N)=0$，称为零测集；某性质$\mu$几乎处处成立，是指不满足性质的点全部能装进某个可测零测集。
但“可以忽略”永远是相对于具体测度而言的：在勒贝格测度下单点$\{0\}$质量为零，可以忽略；在$\delta_0$下它的质量却是$1$，恰好是唯一重要的点。
还要注意例外能够合并的规模：可数个零测集的并仍是零测集，所以可数多个“几乎处处成立”的结论可以共用一个例外范围；不可数个零测集的并则没有这种保证。

另一个相关的性质是完备性。
若每个可测零集的任意子集都仍然可测，就称这个测度空间完备。
不完备时，可能有一个集合包含在零测集里，它本身却不可测——我们知道它“应当质量为零”，却无法在规则内给它赋值。
补救办法是完备化：把这些缺失的零测子集、以及它们与原可测集的并全部补进可测范围，并按外面的零包络赋零。
补进去以后，原来集合的测量值一个都不变，只是可测范围成为满足要求的最小扩展。
实轴上的具体构造前一章已经做过，这里只需理解它的作用：有了完备化，在零测集上随意改动一个函数的值，得到的函数仍然可测；否则“几乎处处相同”的两个函数可能一个可测、一个不可测。
因此使用几乎处处等价的函数时，要么挑一个可测的代表来用，要么明确进入完备化空间，不能简单地以为“反正是零测，怎样都行”。

\subsection{测度与拓扑结构}

仅有可测结构时，我们知道一个集合能否测量，却说不清它是开是闭、是远是近。
后面比较测度时，有些指标（例如全变差）完全不需要位置信息，另一些指标（例如Wasserstein）则必须知道两点之间的距离。
本小节把这两类条件分开：先说明开集怎样进入可测结构，再定义“测度集中在哪里”的\term{支持}。

\subsubsection{Borel结构与标准Borel范围}

\paragraph{从开集生成Borel结构}

当$X$本身带有拓扑（从而能谈开集、闭集）时，最自然的可测范围是让所有开集都可测，再封闭出其余集合。
这个由全体开集生成的$\sigma$代数称为Borel结构：
\[
 \mathcal B(X)=\sigma(\{U\subseteq X:U\text{ 是开集}\}).
\]
于是开集、闭集以及对它们作可数并、可数交、取补得到的集合都可测。
在实轴上，开区间都是开集，这个定义与前一章“由开区间生成$\mathcal B(\mathbb R)$”完全一致。
在离散空间里每个单点本身就是开集，所以Borel结构等于幂集；实轴的Borel结构却不等于全部子集，不可测集依然存在。
请特别注意：选定Borel结构只决定了“哪些集合可测”，它本身\emph{没有指定任何测度}；同一个$\mathcal B(\mathbb R)$上既可以放长度，也可以放正态分布或一个点质量。

选定一个\term{Borel测度}后再做完备化，可测范围还可能扩大。
例如勒贝格可测集族$=$实轴Borel集族$+$所有勒贝格零集的子集，比Borel结构更大。
因此要分清三个层次：拓扑规定开集，Borel$\sigma$代数由开集生成，完备化还要依赖具体选了哪个测度。
后文说“在标准Borel空间上成立”时，指的是中间这一层；不能不加核对地把完备化后才出现的集合代进去。

\paragraph{\term{Polish空间}与\term{标准Borel空间}}

后面的距离比较和测度分解经常需要一个“足够规整、又足够常见”的舞台，为此引入两个条件。
一个拓扑空间若既能找到可数个点稠密地铺在其中（可分，因而开集可以由可数个基本邻域拼出），又能配上一个与拓扑相容的完备度量（柯西列的极限仍落在空间内），就称为Polish空间。
本章用到\term{Wasserstein距离}时，就直接在“给定度量完备且可分”的空间上工作。
这两个条件各管一件事，谁也代替不了谁，更不可能仅仅因为“存在Borel$\sigma$代数”就自动成立。

可测结构之间也能谈“本质相同”。
若双射$h:(X,\mathcal A)\to(Y,\mathcal B)$把可测集映成可测集、且可测集的原像也都可测，也就是每个$B\in\mathcal B$有$h^{-1}(B)\in\mathcal A$、每个$A\in\mathcal A$有$h(A)\in\mathcal B$，则称$h$为\term{可测同构}。
它表示两边的“可测拼图”能一一对应地拼回去，不要求保持距离或其他几何性质。
若一个可测空间能与某个Polish空间的Borel空间可测同构，就称它为标准Borel空间。
实数、欧氏空间、有限或可数离散空间，以及Polish空间的Borel子集都是常见例子。
这里只陈述这一范围并直接使用，不证明其分类定理。
关键直觉是：标准Borel保证“测度上的拼图足够规整”，后面的测度分解等结论可以用；但它不指定距离，也不保证你随手选的距离完备，几何比较仍要另给度量。

\subsubsection{支持与\term{全支撑}}

知道一个测度在哪些集合上有质量后，还想直观地说出“质量集中在空间的哪一片”。
支持就是这个“质量所在地”的精确说法，它还能解释连续函数的几乎处处相等等价于什么。

\paragraph{邻域定义的支持及其全质量}

对可分度量空间上的有限Borel测度$\mu$，定义
\[
 \operatorname{supp}\mu
  =\{x\in X:\text{每个包含 }x\text{ 的开邻域 }U\text{ 都有 }\mu(U)>0\}.
\]
换句话说，支持里的点“走到多近都躲不开质量”；支持外的点则存在一个开邻域完全没有质量。
正因为支持外的点各自被一个零质量开邻域包住，这些邻域的并仍是开集，所以支持是闭集。
支持必须由邻域定义，而不能简单说成“正质量单点的集合”——连续分布下每个单点质量都是零，支持却可以是整个区间。

支持还几乎承载了全部质量。
利用可分性给出的可数拓扑基$\{V_j\}$：对支持外的每个点，先取一个零质量开邻域，再从中选出一个包含该点的基开集。
这样支持外的全部点就被可数个零质量基开集盖住，由可数次可加性，
\[
 \mu(X\setminus\operatorname{supp}\mu)=0.
\]
这一步“先用可数基把任意覆盖变成可数覆盖”不能省略——直接把不可数个零测开集相加是不合法的。
作为对照：均匀分布在$[0,1]$上的测度，每个单点质量都是零，支持却是整个$[0,1]$；$\delta_0$每个邻域只有在含$0$时才有质量，支持就是$\{0\}$。
支持、正质量单点的集合、以及任意一个质量为全重的集合，是三个不同概念，不要混为一谈。

\paragraph{全支撑与连续函数的逐点辨认}

若支持扩到整个空间，$\operatorname{supp}\mu=X$，等价于每个非空开集都有正质量，称$\mu$具有全支撑。
这个条件有一个很实用的推论：在全支撑下，两个连续函数若几乎处处相等，就处处相等。
反证很直接。设连续函数$f,g$几乎处处相等，却在$x_0$处不同。
连续性保证差值$|f-g|$在$x_0$的某个开邻域$U$内仍然严格为正；全支撑又给$\mu(U)>0$。
于是$f\neq g$的地方有正质量，与几乎处处相等矛盾。
这说明：把连续函数看成“只差零测集”的等价类时，全支撑能帮我们找回逐点唯一的代表。

没有全支撑时结论就要打折扣。
在$\delta_0$下，$f(x)=x$与$g(x)=0$几乎处处相等（它们只在$x=0$以外不同，而那里没有质量），却显然不是同一个函数。
一般地，只要在支持上连续且相等，就足以推出几乎处处相等；但要想把结论恢复到整个空间的每一个点，才需要全支撑这类额外条件。

\subsection{一般积分的适用条件}

积分的构造思想在一般测度上并不改变：先在取值简单的函数上按集合质量加权，再用简单函数逼近一般函数。
本小节只回答三个问题：函数满足什么条件才能与可测结构配合、积分什么时候是有限实数、以及前一章的三条积分极限定理在一般空间还成不成立。

\subsubsection{\term{可测映射}与积分条件}

\paragraph{原像判据与可测复合}

函数或映射要参与测量，先要保证它不会把输出端的可测事件拉回输入端的不可测集合。
设$f:(X,\mathcal A)\to(Y,\mathcal B)$，若对每个$B\in\mathcal B$都有
\[
 f^{-1}(B)=\{x\in X:f(x)\in B\}\in\mathcal A,
\]
就称$f$可测。
这里的$f^{-1}(B)$只是“原像集合”的记号，完全不要求$f$有逆函数。

为什么只需在生成族上检查？因为取原像与集合运算可交换：
\[
 f^{-1}(Y\setminus B)=X\setminus f^{-1}(B),\qquad
 f^{-1}\!\left(\bigcup_nB_n\right)=\bigcup_nf^{-1}(B_n).
\]
只要生成$\mathcal B$的那些集合的原像都可测，反复取补和可数并以后，整个$\mathcal B$中集合的原像就都可测。
请注意这与后面的“测度唯一性”是两种推理：前者只靠原像保持集合运算，后者还要靠测度的可加性和有限性。

可测映射复合后仍可测，因为$(g\circ f)^{-1}(C)=f^{-1}(g^{-1}(C))$；连续映射在两边都取Borel结构时自然可测。
对实值函数$f:X\to\mathbb R$，判据还能进一步简化：半直线$\{t:t>a\}$随有理数$a$生成Borel结构，所以只需验证$\{f>a\}\in\mathcal A$对所有有理数$a$成立。
这就把抽象的“可测映射”接回了前一章熟悉的“可测函数”，也为后面计算推前测度和测度核铺平了道路。

\paragraph{非负积分与绝对可积条件}

对非负可测函数，积分仍按“分层称重”定义。
若简单函数$s=\sum_i a_i\mathbf 1_{A_i}$中各层$A_i$可测、两两不交，系数$a_i\geq0$有限，就令
\[
 \int_X s\,d\mu=\sum_i a_i\mu(A_i).
\]
这就是“函数在第$i$层取$a_i$，该层质量为$\mu(A_i)$，乘积相加”。
一般非负函数由简单函数从下逼近，积分取这些逼近值的上确界，允许结果为$+\infty$；二进制逼近和“积分与简单表示无关”的论证与实轴相同，不再重写。

允许函数取负值时，要拆成正负两部分：
\[
 f^+=\max\{f,0\},\qquad f^-=\max\{-f,0\},\qquad f=f^+-f^-.
\]
两部分各自作非负积分，只在它们不都为$+\infty$时才相减。
如果
\[
 \int_X|f|\,d\mu<\infty,
\]
称$f$绝对可积；这时两部分积分都有限，$\int f\,d\mu$才是一个有意义的有限实数，线性和三角不等式等性质也随之成立。
一个常用的充分条件是：在有限测度上，有界可测函数一定可积，因为
\[
 \int_X|f|\,d\mu\leq\|f\|_\infty\,\mu(X)<\infty.
\]
这套积分构造本身并不要求$\sigma$有限。
另外提醒：有正有负的是\emph{被积函数}；\ref{sec:measure-signed-total-variation}小节把测度本身允许取负值，得到的是另一个对象（有符号测度），两者不要混淆。

积分对测度的可数混合也是线性的。
若$\mu=\sum_i a_i\mu_i$且$f\geq0$可测，则
\[
 \int_X f\,d\mu=\sum_i a_i\int_X f\,d\mu_i.
\]
这个等式先对指示函数成立（就是测度定义），再对简单函数由有限可加成立，最后对一般非负函数由单调逼近成立。
对有符号$f$，只有在合成测度下绝对可积时才能这样拆开。
取点质量混合时，它退化成最直观的加权求和：
\[
 \int_X f\,d\!\left(\sum_i w_i\delta_{x_i}\right)=\sum_i w_i f(x_i).
\]
一般空间上的积分因此统一了“连续铺开”的累计和“在离散点上”的加权平均。

\subsubsection{极限操作的可测性与积分交换}

使用函数列时要把两件事分开：先问极限函数本身能不能测，再问极限符号能不能搬进积分号。
两件事各有各的条件，“每一项都可测、都可积”并不自动保证它们的任意汇总合法。

\paragraph{可数极限与不可数上确界的边界}

对可数个可测函数，阈值集合的可数并、可数交仍然可测，所以
\[
 \sup_n f_n,\quad\inf_n f_n,\quad\limsup_n f_n,\quad\liminf_n f_n
\]
都是可测函数；逐点极限存在时，极限函数也可测。
非负级数的部分和是一种特殊的可数极限，因此前面“可数混合的积分可以逐项求和”有合法的函数对象作保障。

但$\sigma$代数只对\emph{可数}次运算封闭，不可数个可测函数的逐点上确界未必可测。
实轴Borel结构上有个很说明问题的反例：取一个非Borel集合$N$，每个单点示性函数$\mathbf 1_{\{x\}}$都是可测的，但
\[
 \sup_{x\in N}\mathbf 1_{\{x\}}=\mathbf 1_N
\]
却是不可测函数。
这划出了一般规则的边界：不可数次优化或取极限要保证可测性，必须在具体场景中另加可分性、可测选择等条件，不能因为“每一项都可测”就默认其不可数汇总也可测。

\paragraph{三条积分极限规则的适用条件}

前一章的三条积分极限定理，其证明只用到简单函数逼近、测度的可数可加和可数极限可测，不依赖实轴的几何，因此在任意一般测度空间原样成立。
这里不重复证明，只把使用条件强调清楚。

若非负可测函数列单调上升$f_n\uparrow f$，单调收敛定理给出
\[
 \int f_n\,d\mu\uparrow\int f\,d\mu,
\]
两边一起取$+\infty$也允许。
对任意非负可测函数列，Fatou不等式
\[
 \int\liminf_n f_n\,d\mu\leq\liminf_n\int f_n\,d\mu
\]
始终成立，它只给“先取极限”一侧的下界，甚至不要求$f_n$真的有逐点极限。
支配收敛定理则要求三个条件同时具备：$f_n$几乎处处收敛到可测的$f$、存在\emph{同一个}非负可积函数$g$、且$|f_n|\leq g$几乎处处成立；这时积分可以取有限极限，并且
\[
 \int|f_n-f|\,d\mu\to0.
\]
证明时可以把可数个“不收敛、不被支配、取无穷值”的例外合并成一个零测集，在其外使用Fatou式论证，从而避免无穷相减。
请特别注意“同一个支配函数”这个要求不能由“每个$f_n$各自可积”替代。
本小节始终固定测度、只改变被积函数；至于“测度本身变化时积分怎样收敛”，要等到\ref{sec:measure-sequence-convergence}节讨论。

\subsection{乘积结构与测度唯一性}

后面比较测度时，经常先在“矩形”这类容易计算的集合上构造或验证一个测度，再断言它在整个空间上被唯一确定；讨论两个空间的配对时，也需要联合空间的可测结构。
本小节建立三件共同工具：简单集合上的相等怎样推广到全部集合、两个测度空间怎样相乘、以及更多坐标下的柱集接口。
要先强调一点：给定两个边缘测度，乘积测度只是联合测度的一种，耦合和测度核会给出许多别的联合测度。

\subsubsection{生成族上的测度唯一性}

\paragraph{有限测度的$\pi$系统唯一性}

在半直线、矩形这类简单集合上赋值很容易，但我们真正想确定的是整个$\sigma$代数上的测度。
自然的问题是：如果在一族简单集合上两项测度取值相同，能否断定它们处处相同？
答案在一定条件下是肯定的。
若集合族$\mathcal C$对有限交封闭（$C_1,C_2\in\mathcal C\Rightarrow C_1\cap C_2\in\mathcal C$），称它为$\pi$系统；半直线和可测矩形都满足这个性质。

\begin{theorem}[$\pi$系统上的测度唯一性]
\label{thm:measure-pi-lambda-uniqueness}
设$\mathcal C$是包含$X$且生成$\mathcal A$的$\pi$系统，$\mu,\nu$是$(X,\mathcal A)$上的有限测度。
若$\mu(C)=\nu(C)$对每个$C\in\mathcal C$成立，则$\mu=\nu$在整个$\mathcal A$上成立。
\end{theorem}

“包含$X$”这一条不能漏，它保证两项测度的总质量相等。
证明的核心思路值得了解，因为它在乘积和核的构造中会反复用到。
令
\[
 \mathcal D=\{A\in\mathcal A:\mu(A)=\nu(A)\}
\]
为所有“两项测度取值相同”的集合。
因为总质量有限且相等，$A\in\mathcal D$时补集也相等：
\[
 \mu(X\setminus A)=\mu(X)-\mu(A)=\nu(X)-\nu(A)=\nu(X\setminus A).
\]
又因为测度可数可加，两两不交集合并起来后仍然相等，加上$X\in\mathcal D$，$\mathcal D$就是一个对补集和不交可数并封闭的$\lambda$系统。
$\pi$—$\lambda$定理说，包含$\pi$系统$\mathcal C$的最小$\lambda$系统恰好是$\sigma(\mathcal C)=\mathcal A$，所以全部可测集都在$\mathcal D$中。
直观上讲：在有限总质量下，只要在“积木”上对齐，整体质量、补集和不交拼接就全部对齐了。

$\sigma$有限版本要用，但必须写清共同的有限覆盖：存在$E_n\in\mathcal C$、$E_n\uparrow X$，使得$\mu(E_n)=\nu(E_n)<\infty$。
做法是先在每个$E_n$内用有限版本比较截断测度$\mu(\cdot\cap E_n)$与$\nu(\cdot\cap E_n)$，再令$n\to\infty$拼出全空间。
只知道两项测度各自$\sigma$有限、却没有同一套有限盖子，并不能保证它们在生成族上相等。

\subsubsection{有限乘积与累次积分}

\paragraph{乘积测度与一般Tonelli/Fubini接口}

要同时描述两个空间上的配对，先规定联合空间里哪些集合可测。
设$(X,\mathcal A,\mu)$与$(Y,\mathcal B,\nu)$都是$\sigma$有限测度空间。
可测矩形$A\times B$生成乘积$\sigma$代数$\mathcal A\otimes\mathcal B$；在矩形上自然定义
\[
 (\mu\otimes\nu)(A\times B)=\mu(A)\nu(B),
\]
也就是“两块质量相乘”，其中任一因子质量为零时乘积记为零。
在$\sigma$有限条件下，这个矩形赋值能唯一扩展成$\mathcal A\otimes\mathcal B$上的乘积测度$\mu\otimes\nu$。
唯一性来自矩形族是$\pi$系统、且两个因子的有限覆盖$X_k,Y_l$给出联合空间的有限矩形覆盖$X_k\times Y_l$；存在性是矩形赋权的标准扩展，前一章在欧氏空间已经构造过，这里直接引用。
但请记住：乘积测度只是“两边独立地各算各的”这一种联合方式；同样的边缘完全可以对应许多非乘积的联合测度。

当$X,Y$都是可分度量空间时，拓扑Borel结构和乘积可测结构是一致的：
\[
 \mathcal B(X\times Y)=\mathcal B(X)\otimes\mathcal B(Y).
\]
一方面坐标投影连续，开集的原像可测，所以Borel矩形都在乘积拓扑的Borel结构中；另一方面，两个可数拓扑基的矩形构成乘积拓扑的可数基，任何开集都是可数个开矩形之并。
有了这个等式，后面在联合空间上谈Borel耦合、下半连续成本和弱极限时，用的都是同一套可测集合，不会出现“事件可测性对不上”的问题。

积分换序的定理直接沿用前一章，这里只确认一般空间下的条件。
对非负的乘积可测函数$f$，Tonelli定理保证：每个截面可测，内层积分作为另一坐标的函数可测，而且两种累次积分与联合积分相等，
\[
 \int_{X\times Y}f\,d(\mu\otimes\nu)
 =\int_X\!\left(\int_Y f(x,y)\,d\nu(y)\right)d\mu(x)
 =\int_Y\!\left(\int_X f(x,y)\,d\mu(x)\right)d\nu(y),
\]
即使积分值为$+\infty$也成立——这正是非负性带来的便利，不用担心正负无穷相消。
对有符号函数，则必须先用
\[
 \int_{X\times Y}|f|\,d(\mu\otimes\nu)<\infty
\]
保证绝对可积，Fubini定理才给出有限的换序结论；而且截面可积只对几乎处处的剩余坐标成立。
特别要避免一个常见错误：只看到某一种次序下的有符号累次积分存在，不足以换序，必须先确认联合绝对积分有限。

\paragraph{乘积完备化中的截面解释}

如果$f$只在乘积测度的完备化中可测，处理截面时要小心。
可以先选一个与$f$乘积几乎处处相同、但对原乘积$\sigma$代数可测的代表$\widetilde f$，对$\widetilde f$使用上述定理。
两者的差异集包含在某个原乘积可测零集$N$内；对$\mathbf 1_N$用Tonelli可知，$N$的截面在几乎处处的剩余坐标下可测且质量为零。
这一步给出了处理$f$截面所必需的零包络，不能直接假定差异集的截面在原来的因子结构中可测。
在这些“好”坐标上，$f$与$\widetilde f$的截面只在零截面的子集上不同；必要时还要进入因子测度的完备化，才能把这个子集当作可测零集。
所以不能宣称每个截面都合法，也不能由“两个因子分别完备”推出乘积$\sigma$代数本身已经完备。
把这层细微区别交代清楚以后，后面只在条件已经核对过的场合调用换序。

\subsubsection{多坐标与柱集的接口}

两个空间的乘积可以重复进行，得到任意有限个空间的乘积；但要描述无穷多个坐标（例如一条随时间展开的轨迹），得先说清可测范围和不同维数之间怎样保持一致。
这里只建立接口，不提前讲随机过程的扩张理论。

\paragraph{有限坐标信息与相容边缘}

在乘积空间$\prod_{i\in I}X_i$上，只约束有限多个坐标的集合称为柱集；全体柱集生成这里的乘积$\sigma$代数。
每个坐标投影都可测，只依赖有限个坐标的函数也仍可用原像判据检验可测性。
设想对每个有限指标集合$J$都指定一项联合测度$\mu_J$。
它们不能各自任意指定：从较大的有限坐标集投影回$J$以后，必须恰好恢复$\mu_J$，这就是相容性。
它保证“在更粗的有限观察下看到的边缘分布”，在不同维数的赋值中保持一致。

有限维相容是构造无穷坐标测度时必须通过的检查，但仅靠它不能在任意可测空间上断言扩张存在；还需要标准Borel这类规整条件，扩张定理、路径的正则性以及随机过程解释都留给后文的概率论。
到此为止，后面比较与映射所需的可测范围、一般赋权、积分条件和联合计算工具已经备齐；下一节起，我们固定同一个可测空间，开始比较上面不同测度之间的关系。

\section{同一空间上的测度比较}
\label{sec:measure-same-space-comparison}

本节固定共同底集与可测结构$(X,\mathcal A)$。
先判断两项测度的零集是否相容，再研究一项测度的质量有多少能由相对密度表示；随后定量比较整体差异：测度之差与全变差累计正负赋权差，\term{Hellinger距离}读取平方根密度差，\term{KL散度}读取有向对数比，\term{Wasserstein距离}再额外使用共同底层距离。
最后用允许的测试函数族综合这些工具，说明它们各自能控制哪些积分读数。
这些工具按所问的问题分工：关系、表示和综合回看并不是互斥的测度类别。

零集关系本身对一般非负测度有意义；\term{Radon–Nikodym密度}和\term{Lebesgue分解}采用非负$\sigma$有限版本。
为解释两项测度之差，本小节再补有限\term{有符号测度}；其后的距离与散度统一以概率为对象。
共同空间只保证可以比较同一个事件，既不自动给出密度，也不自动提供可用的几何距离。

\subsection{\term{绝对连续}、等价与\term{相互奇异}}

本小节比较可测零集与质量集中范围。
两个子项分别处理零集的单向或双向蕴含，以及质量能否集中在互补的可测范围内：绝对连续有方向，等价要求双向共享零集，相互奇异要求可测集中范围分离。
这些关系可以分别检验，不组成互斥且穷尽的分类。

\subsubsection{绝对连续与\term{等价测度}}

\paragraph{绝对连续的定义与方向}

设$\mu,\nu$是同一可测空间上的非负测度。
先问：若$\mu$把某个可测集合视为零质量，$\nu$是否也必须如此？
若对每个$A\in\mathcal A$都有
\[
 \mu(A)=0\ \Longrightarrow\ \nu(A)=0,
\]
称$\nu$相对于$\mu$绝对连续，记为$\nu\ll\mu$。
方向写在“$\nu$相对于$\mu$”这一关系中；比较的是整个可测集合的零质量，而不是逐点密度值的大小。

若已经存在非负可测函数$f$使
\[
 \nu(A)=\int_A f\,d\mu\qquad(A\in\mathcal A),
\]
则在$\mu(A)=0$的集合上积分为零，立刻得到$\nu(A)=0$，所以这种密度表示足以推出$\nu\ll\mu$。
这只是一个方向；能否从零集蕴含反推出这样的密度函数，需要后面的Radon–Nikodym定理及其假设。

绝对连续并非连续空间专有。
实轴上相对于勒贝格测度有密度的权重满足绝对连续，\term{点质量}$\delta_0$却不满足：$\lambda(\{0\})=0$而$\delta_0(\{0\})=1$。
在有限或可数离散底集上取计数测度作参考，也可以同样定义绝对连续。
还要注意只检验单点并不够。
在实轴上令$P,Q$分别为均匀分布于$[0,1]$与$[2,3]$的概率，所有单点在两者下质量都为零，但$Q([0,1])=0$而$P([0,1])=1$，所以$P$不对$Q$绝对连续。
一般可测集合可能不可数，不能用只对可数并成立的可加性，把它的质量改写成全部单点质量之和；单点检验没有覆盖定义中的全体集合量词。

\paragraph{等价测度与共同零集}

在单向绝对连续的基础上，若$\mu\ll\nu$与$\nu\ll\mu$同时成立，记为$\mu\sim\nu$，称两测度等价；这恰好表示两者具有相同的可测零集。
等价并不意味着相等：非零集合的质量可以按不同比例分配，只是“谁质量为零”保持一致。

若已有合法表示$\nu=f\mu$且$f>0$在$\mu$几乎处处成立，则可补足反向绝对连续。
设$\nu(A)=0$，即$\int_Af\,d\mu=0$。对每个正整数$m$，由
\[
 f\,\mathbf 1_A\geq\frac1m\,\mathbf 1_{A\cap\{f\geq1/m\}}
\]
与积分单调性得到
\[
 \mu\bigl(A\cap\{f\geq1/m\}\bigr)\leq m\int_Af\,d\mu=0.
\]
正值层$\{f>0\}=\bigcup_m\{f\geq1/m\}$上$A$的部分是可数个零测集的并，质量仍为零；与$f>0$几乎处处成立合并，得到$\mu(A)=0$。
所以用几乎处处为正的密度重新加权，只改变质量大小、不改变零集。
这是“已经有密度且几乎处处为正”的充分判据；密度何时存在、何时能取倒数，将在相对密度小节补齐条件。

\subsubsection{相互奇异}

\paragraph{相互奇异与可测集中范围}

与共享零集相对的另一个问题是：两项质量能否分别集中在互补范围。
若存在可测集合$E$使
\[
 \mu(X\setminus E)=0,\qquad \nu(E)=0,
\]
称$\mu$与$\nu$相互奇异，记为$\mu\perp\nu$。
前一式表示$\mu$的全部质量集中在$E$上，后一式表示$\nu$的全部质量集中在补集$X\setminus E$上。

对勒贝格测度$\lambda$与$\delta_0$，取$E=\mathbb R\setminus\{0\}$即可：$\lambda(\{0\})=0$，$\delta_0(\mathbb R\setminus\{0\})=0$。
这里$E=\mathbb R\setminus\{0\}$与$\{0\}$不相交，但$E$的闭包是整条实轴，与$\{0\}$的闭包仍有交点。
相互奇异要求的是可测集中范围分离，并不要求两个范围的闭包不相交。

一项测度还可能同时含有相对于参考绝对连续的质量和与参考相互奇异的质量，这时它整体既不绝对连续也不相互奇异。
奇异性也不要求奇异部分由单点组成：后面的Lebesgue分解处理的是整个测度的两部分，奇异部分可以是支撑在零测的不可数集合上的连续质量。

\subsection{相对密度与测度分解}

本小节回答四个问题：密度怎样取得，密度值表示什么，怎样用它改写积分，以及哪些质量必须另外分解。
采用非负$\sigma$\term{有限测度}版本的Radon–Nikodym定理：存在结论直接陈述，不展开完整存在证明；几乎处处唯一性、单点解释和换参考公式分别说明依据。
这里的“导数”是相对积分权重，不是普通坐标导数。

\subsubsection{相对密度的存在与几乎处处唯一}

\paragraph{相对密度的存在与有限性}

绝对连续$\nu\ll\mu$只描述零集关系。
更进一步，希望用一个函数恢复$\nu$在所有可测集合上的权重。

\begin{theorem}[Radon–Nikodym定理]
\label{thm:measure-radon-nikodym}
设$\mu,\nu$是同一可测空间上的非负$\sigma$有限测度，且$\nu\ll\mu$。
则存在非负可测函数$f:X\to[0,\infty]$，使
\[
 \nu(A)=\int_A f\,d\mu\qquad(A\in\mathcal A).
\]
\end{theorem}

这个$f$记为$f=\dfrac{d\nu}{d\mu}$，称为$\nu$相对于$\mu$的Radon–Nikodym密度。
$\sigma$有限是这里所采用定理的充分假设，不写成一切密度表示的必要条件；对已有连续密度或可数离散权重的情形，可直接把恢复集合权重的已知函数辨认为Radon–Nikodym密度。

目标测度$\nu$的$\sigma$有限性还能保证密度几乎处处有限。
取可测$E_k$覆盖$X$且$\nu(E_k)<\infty$，令$N=\{f=\infty\}$。
若$\mu(N\cap E_k)>0$，则
\[
 \nu(E_k)=\int_{E_k}f\,d\mu\geq\int_{N\cap E_k}f\,d\mu=+\infty,
\]
与$\nu(E_k)<\infty$矛盾，故每个$\mu(N\cap E_k)=0$。
由可数覆盖$\mu(N)=0$，可以选取$f$在$\mu$几乎处处有限的代表，这也为后面等价测度的倒数公式准备了条件。

\paragraph{密度的几乎处处唯一}

若$f,g$都恢复$\nu$在每个可测集合上的质量，则它们只能在零集上不同。
除去二者取无穷值的共同$\mu$零集，取$\mu$有限的可测覆盖$\{M_j\}$和$\nu$有限的可测覆盖$\{N_k\}$，在每个$M_j\cap N_k$内比较；再按$f,g\leq l$的有界取值层和$f\geq g+1/m$的间隔层细分。
这些限制保证比较时积分有限，并把“严格大于”转化为至少$1/m$的正间隔。

对这样的可测层$D$，两种表示给出$\int_Df\,d\mu=\nu(D)=\int_Dg\,d\mu$。
若$\mu(D)>0$，正间隔迫使左侧比右侧至少多$\mu(D)/m$，产生矛盾；所以每个这样的层都是零测集。
可数合并各质量块、取值层和间隔层，得到$\mu(\{f>g\})=0$；交换$f,g$再得$\mu(\{g>f\})=0$。
因此
\[
 \frac{d\nu}{d\mu}=f=g\quad\mu\text{几乎处处}.
\]
在$\mu$零集上改变代表值不会改变任何集合权重；导数记号并没有指定每个点的数值，下一目就用单点质量把这一区别算清。

\subsubsection{密度值与单点质量}

\paragraph{密度系数与单点质量的区别}

固定$\nu=f\mu$并先确认$\{x\}$可测，则
\[
 \nu(\{x\})=\int_{\{x\}}f\,d\mu.
\]
这个公式比较的是单点集合的累计质量与恢复质量的函数系数，两者一般不是同一个数。

在$[0,1]$上以勒贝格测度为参考，取概率密度$p(x)=1$和$q(x)=2x$。
两项测度给每个单点都赋零质量，但在区间$[0,1/2]$上的质量分别为
\[
 \int_0^{1/2}1\,dx=\frac12,\qquad
 \int_0^{1/2}2x\,dx=\frac14.
\]
密度表示的是权重如何沿集合分布，并不是“在该点的概率”。

改换参考条件后结论才不同。
若$0<\mu(\{x\})<\infty$，单点上的积分退化为$f(x)\mu(\{x\})$，此时存在密度表示才可由
\[
 f(x)=\frac{\nu(\{x\})}{\mu(\{x\})}
\]
决定点系数，这正是有限或可数离散权重中的质量比。
若$\mu(\{x\})=0$，改变$f(x)$并不改变测度，几乎处处唯一性不会指定该点的值。

还可用$\delta_0$相对于勒贝格测度检验表示是否可能。
任何普通非负勒贝格密度$f$在$\{0\}$上的积分都为零，却应恢复质量$\delta_0(\{0\})=1$，所以不存在这样的普通可测密度。
广义函数理论中的Dirac记号需要另建分布框架，不能在这里冒充Radon–Nikodym可测函数。
缺少整个测度的参考密度，也不能靠在某个点填写无穷值补救；不能表示的奇异质量更未必只由单点构成。

\subsubsection{积分的换测度公式}

\paragraph{用相对密度改写积分}

取$f=d\nu/d\mu$。
先在指标函数$g=\mathbf 1_A$上由Radon–Nikodym表示直接得到
\[
 \int_X g\,d\nu=\nu(A)=\int_A f\,d\mu=\int_X gf\,d\mu.
\]
这是固定底层可测空间、把同一个被积对象改由另一参考测度累计，没有移动输入点，也没有重新生成输出集合。

对非负简单函数把指标恒等式有限相加；再对非负可测$g$取简单函数单调逼近，两侧分别使用单调收敛，得到
\[
 \int_X g\,d\nu=\int_X gf\,d\mu,
\]
允许两边共同为$+\infty$。
可能出现的$0\cdot\infty$按非负积分的零贡献约定处理，不在某个零密度点上强行使用普通有限数乘法。
对有符号$g$则分别计算正负部分，只在不出现$+\infty-(+\infty)$的合法范围内合并。
特别地，
\[
 \int|g|\,d\nu=\int|g|f\,d\mu
\]
使绝对可积条件在两种写法中对齐；可测性、存在密度和积分有限仍是三项不同检查。

\paragraph{重加权、限制与归一化}

换参考公式也可以反向使用：用非负可测权重$w$在$\mu$上构造新测度
\[
 \nu(A)=\int_A w\,d\mu.
\]
指标上的非负可加性与单调收敛保证$\nu$可数可加。
若归一化常数
\[
 Z=\int_Xw\,d\mu
\]
满足$0<Z<\infty$，则$Q(A)=\nu(A)/Z$是\term{概率测度}，相对于$\mu$的密度为$w/Z$；$Z=0$或$Z=+\infty$时不能按这个公式归一化，权重可测不保证$Z$有限正。

取$w=\mathbf 1_E$得到\term{限制测度}
\[
 \mu_E(A)=\mu(A\cap E);
\]
当$0<\mu(E)<\infty$时，归一化后只保留$E$内的相对权重。
底集和可测事件都没有变，被改变的是赋权；它与把输入点送到另一空间的推前不同，也与只把测度限制到较粗子$\sigma$代数不同。
在$[0,1]$上的均匀概率下取$w=2\mathbf 1_{[0,1/2]}$，则$Z=1$，新概率在$[0,1/2]$上质量为$1$、在$(1/2,1]$上质量为$0$；它对原参考绝对连续，却不再等价，因为原来正质量的区域被删去了。

\paragraph{链式换参考与倒数公式}

若$\rho\ll\nu\ll\mu$且三者都在当前非负$\sigma$有限范围内，依次对集合指标应用两次换参考公式，得到
\[
 \frac{d\rho}{d\mu}
  =\frac{d\rho}{d\nu}\frac{d\nu}{d\mu}
  \qquad\mu\text{几乎处处}.
\]
在$d\nu/d\mu=0$处把乘积的代表取为零，避免把未定的$0\cdot\infty$当成通常乘法。

若进一步$\mu\sim\nu$，令$A=\{f=0\}$，由$\nu(A)=\int_Af\,d\mu=0$与反向绝对连续得$\mu(A)=0$，故$f=d\nu/d\mu>0$在$\mu$几乎处处成立；结合前面由$\nu$的$\sigma$有限性得到的$f<\infty$，存在一个共同满测度区域使$0<f<\infty$。
在该区域定义$1/f$，对可测$B$有
\[
 \int_B\frac1f\,d\nu=\int_B\frac1f\,f\,d\mu=\mu(B),
\]
故$d\mu/d\nu=1/f$在$\nu$几乎处处成立；零集上的代表另行约定，不贡献积分。
单向绝对连续不足以取得这一步。
用简单常数例子核对方向：同一有限参考$\mu$上取$\nu=2\mu$、$\rho=3\nu$，则$d\nu/d\mu=2$、$d\rho/d\nu=3$、$d\rho/d\mu=6$，反向$d\mu/d\nu=1/2$。
公式的依据始终是恢复所有可测集合上的积分，不能形式地把微分符号约去。

\subsubsection{Lebesgue分解}

\paragraph{绝对连续部分与奇异部分}

整个测度相对于固定参考未必绝对连续，因此除了表示能写成密度的部分，还要拆出不能这样表示的部分。

\begin{theorem}[Lebesgue分解]
\label{thm:measure-lebesgue-decomposition}
设$\mu,\nu$是同一可测空间上的非负$\sigma$有限测度。
则存在唯一的非负测度$\nu_{\mathrm{ac}}$与$\nu_{\mathrm s}$，使
\[
 \nu=\nu_{\mathrm{ac}}+\nu_{\mathrm s},\qquad
 \nu_{\mathrm{ac}}\ll\mu,\qquad
 \nu_{\mathrm s}\perp\mu.
\]
\end{theorem}

对绝对连续部分应用Radon–Nikodym定理，可写成$\nu_{\mathrm{ac}}=f\mu$，于是每个可测集合满足
\[
 \nu(A)=\int_Af\,d\mu+\nu_{\mathrm s}(A).
\]
两项定理的职责要分开：分解先把测度拆成两部分，Radon–Nikodym只表示其中的绝对连续部分。
若$\nu_{\mathrm s}\neq0$，不能把$f$称为整个$\nu$相对于$\mu$的密度。
这里的唯一性针对两项测度；存在性按当前范围陈述，不展开完整证明。

\paragraph{分解的示例与参考测度依赖}

在$[0,1]$上取参考$\mu=\lambda|_{[0,1]}$及
\[
 \nu=\lambda|_{[0,1]}+\delta_0.
\]
长度部分对$\mu$绝对连续；点质量部分以$\{0\}$与$\mu$相互奇异，由此可辨认
\[
 \nu_{\mathrm{ac}}=\lambda|_{[0,1]},\qquad
 \nu_{\mathrm s}=\delta_0.
\]
集合$\{0\}$总质量为$1$，不含$0$的区间$(a,b)\subset(0,1)$总质量为$b-a$；含$0$的集合要同时累计两部分。
只有两部分相加才能恢复$\nu$对全部可测集合的权重。

分解依赖参考。保持$\nu$不变而改取$\mu'=\nu$，则$\nu$相对于$\mu'$的密度恒为$1$、奇异部分为零。
同一测度在两次分解中的分量不同，说明分解不是脱离参考的固有标签。
奇异性只要求可测集中范围分离，不要求奇异部分由点质量组成，不能把这个Dirac特例的形状推广为全部奇异测度。
下面保持共同可测空间，转向两项概率的整体差异。

\subsection{有限有符号测度与全变差}

要比较同一事件上的质量差，可以先把两项非负赋权相减。
差值允许为负，需要区分原来的非负测度、两测度之差以及它的绝对变化总量；随后再把这个通用对象用于概率之间的\term{全变差距离}。

\subsubsection{有限有符号测度与\term{总变差测度}}
\label{sec:measure-signed-total-variation}

\paragraph{集合赋权的正负分量}

对有限非负测度$\mu,\nu$，令
\[
 \eta(A)=\mu(A)-\nu(A).
\]
沿相同的不交可数分解，$\eta$仍然可数可加，但单个集合上的结果可以为负；这与“函数取正负值而积分测度仍非负”是两种不同的对象变化。
本章把取有限实值且可数可加的集合赋权称为有限有符号测度，不涉及无穷正负质量相减。

\begin{theorem}[\term{Hahn分解}与\term{Jordan分解}]
\label{thm:measure-hahn-jordan}
对有限有符号测度$\eta$，存在可测集合$E_+,E_-$满足$E_+\cup E_-=X$、$E_+\cap E_-=\varnothing$，使$\eta$在$E_+$的每个可测子集上非负、在$E_-$的每个可测子集上非正。
令
\[
 \eta^+(A)=\eta(A\cap E_+),\qquad
 \eta^-(A)=-\eta(A\cap E_-),
\]
则$\eta^+,\eta^-$是相互奇异的非负有限测度，且
\[
 \eta=\eta^+-\eta^-,
\]
这两个Jordan分量作为测度唯一。
\end{theorem}

区域$E_+,E_-$本身的选择不必唯一，但分量测度唯一。
在实轴上取$\eta=\delta_0-\delta_1$，可逐点检查$\eta(\{0\})=1$、$\eta(\{1\})=-1$、$\eta(\mathbb R)=0$，正负分量分别是$\delta_0$与$\delta_1$；全空间质量为零不代表这个集合赋权处处为零。
Jordan分解按符号分，Lebesgue分解相对于另一参考测度按绝对连续与奇异分，两者回答不同问题。

\paragraph{总变差测度与有符号积分}

定义总变差测度与\term{总变差范数}：
\[
 |\eta|=\eta^++\eta^-,\qquad
 \|\eta\|_{\mathrm{var}}=|\eta|(X).
\]
对任意可测集合$A$的有限可测分割$\{A_i\}$，都有
\[
 \sum_i|\eta(A_i)|\leq\eta^+(A)+\eta^-(A)=|\eta|(A);
\]
用$A\cap E_+$与$A\cap E_-$这两个集合分割恰好取得等号，因此$|\eta|(A)$也等于所有分割的绝对赋权和上确界。
它累计正负变化，不让两者抵消。

对满足$\int|g|\,d|\eta|<\infty$的可测$g$，定义
\[
 \int g\,d\eta=\int g\,d\eta^+-\int g\,d\eta^-,
\]
由两部分分别估计得到
\[
 \left|\int g\,d\eta\right|\leq\int|g|\,d|\eta|.
\]
有界可测函数因此满足
\[
 \left|\int g\,d\eta\right|\leq\|g\|_\infty\|\eta\|_{\mathrm{var}};
\]
控制积分要靠总变差，不能只看可能正负抵消的$\eta(X)$。

若$f$可测且$\int|f|\,d\mu<\infty$，令$\eta(A)=\int_Af\,d\mu$，则$f^+\mu$与$f^-\mu$正是它的Jordan分量，故
\[
 |\eta|=|f|\mu.
\]
这把一般有符号赋权接回已建立的非负积分。
对有限有符号测度，“所有参考零集上的赋权为零”等价于其总变差在这些零集上为零，因为可把每个参考零集再按Hahn正负区域分开检查。
本章不引入一般的向量测度理论。

\subsubsection{全变差距离与测试积分}

现在限定为同一可测空间上的概率$P,Q$。
取共同有限参考$\lambda=P+Q$，由Radon–Nikodym定理写
\[
 p=\frac{dP}{d\lambda},\qquad q=\frac{dQ}{d\lambda}.
\]
$P+Q$有限，$P,Q$都对它绝对连续，密度存在；测度之差的密度为$p-q$，总质量抵消会给出概率全变差特有的二分之一系数。

\paragraph{事件质量差与密度公式}

定义
\[
 \operatorname{TV}(P,Q)=\sup_{A\in\mathcal A}|P(A)-Q(A)|.
\]
候选是同一可测空间中的全部事件，比较的是同一事件在两种权重下的差。
由$0\leq P(A),Q(A)\leq1$可得$0\leq\operatorname{TV}\leq1$。
三角不等式由实数三角不等式逐项使用再取上确界得到；距离为零意味着所有事件质量相同，即$P=Q$；对称性由绝对值取得。
所以TV确实构成概率测度之间的距离。

在共同参考下令$h=p-q$。
总质量都为$1$给出$\int h\,d\lambda=0$，所以正部积分与负部积分相等。
取$A_+=\{p\geq q\}$，任意事件上的积分差都不超过正部总量，而$A_+$恰好达到该量。
正负两部分相加得到
\[
 \operatorname{TV}(P,Q)
  =\frac12\int|p-q|\,d\lambda
  =\frac12|P-Q|(X).
\]
二分之一来自等总质量；换另一个共同有限参考，事件差上确界不变，所以数值与参考选择无关。
推广到总质量不相等的有限测度时必须重新声明归一化约定，不能直接沿用这个$1/2$。

\paragraph{有界观测量的误差控制}

对取值于$[a,b]$的可测函数$g$，利用两项概率总质量相同，先减去常数$(a+b)/2$，再由共同参考下的密度差估计，可得
\[
 \left|\int g\,dP-\int g\,dQ\right|
  \leq(b-a)\operatorname{TV}(P,Q).
\]
特别地，$0\leq g\leq1$时系数为$1$；令$a=-\|g\|_\infty$、$b=\|g\|_\infty$，恢复一般有界函数的$2\|g\|_\infty$界。
无界观测还需要额外条件。

因为所有$0\leq g\leq1$的积分差都受TV控制，而正差集合的指标$g=\mathbf 1_{A_+}$达到上确界，所以
\[
 \operatorname{TV}(P,Q)=\sup_{0\leq g\leq1}\left|\int g\,dP-\int g\,dQ\right|.
\]
用$g=2\mathbf 1_{A_+}-1$以及有界函数界，还有
\[
 \sup_{\|g\|_\infty\leq1}\left|\int g\,dP-\int g\,dQ\right|
  =2\operatorname{TV}(P,Q).
\]
两种测试类的归一化不同，这一系数在后面比较\term{测度核}作用时还会用到。
在实轴Borel空间上，$\delta_a$与自身距离为零；$a\neq b$时事件$\{a\}$把两点区分，所以$\operatorname{TV}(\delta_a,\delta_b)=1$。
全变差辨认事件质量差异，却不读取$a,b$之间的几何间隔，这为后面的运输比较留下了明确问题。

\subsection{Hellinger距离与KL散度}

仍以同空间概率及共同有限参考$\lambda=P+Q$为对象，沿用$p=dP/d\lambda$、$q=dQ/d\lambda$。
与累计绝对密度差的全变差相比，下面分别使用平方根密度差和有向对数比，建立指标后再比较它们与全变差的关系。
这样不必在每项中重复共同参考的定义，也不把“基于密度”误写成必须有勒贝格密度。

\subsubsection{Hellinger距离}

\paragraph{平方根密度与参考无关性}

约定
\[
 H^2(P,Q)=\frac12\int(\sqrt p-\sqrt q)^2\,d\lambda,
\]
令$H(P,Q)=\sqrt{H^2(P,Q)}$。
利用$\int p\,d\lambda=\int q\,d\lambda=1$展开：
\[
 H^2(P,Q)
  =\frac12\!\left(\int p+\int q-2\int\sqrt{pq}\right)
  =1-\int\sqrt{pq}\,d\lambda.
\]
由积分形式的Cauchy–Schwarz不等式，$0\leq\int\sqrt{pq}\,d\lambda\leq1$，所以$0\leq H\leq1$。

换共同参考不改变$H$。
若改用另一共同$\sigma$有限参考$\eta$，把$\lambda,\eta$同置于$\kappa=\lambda+\eta$下，用Radon–Nikodym链式关系写出两概率相对于$\kappa$的密度；平方根差提出共同因子后再换参考积分，得到同一个$H^2$。
零密度区按积分零贡献处理，不能只凭两个公式外形相似就宣称不变。

\paragraph{距离性质与端点情形}

非负性和对称性直接由公式得到。
$H=0$意味着$\sqrt p=\sqrt q$几乎处处成立，即两密度相同，恢复的全部集合权重相同。
证明三角不等式时引入第三项概率$R$，三者统一以$P+Q+R$为参考，避免在不同参考中直接相加平方根密度；再把两段平方根差相加，对平方积分中的交叉项应用Cauchy–Schwarz即可。
所需积分不等式可由简单函数的有限和情形经逼近建立，不预先调用后章的$L^2$空间理论。

端点值可用极端情形核对：$P=Q$时$H=0$；$P\perp Q$时共同参考下两者的正密度区域不相交，$\sqrt{pq}=0$，故$H=1$。
要注意$H$和$H^2$的区别：本节称为距离的是开平方后的$H$；它与全变差、KL的关系留到\ref{sec:measure-metric-relations}小节。

\subsubsection{KL散度与有限性}

\paragraph{有向定义与积分合法性}

若$P\ll Q$，定义
\[
 D_{\mathrm{KL}}(P\|Q)=\int\log\frac{dP}{dQ}\,dP,
\]
使用自然对数；若$P$不对$Q$绝对连续，约定$D_{\mathrm{KL}}(P\|Q)=+\infty$。
在共同参考下可写成
\[
 D_{\mathrm{KL}}(P\|Q)=\int p\log\frac pq\,d\lambda,
\]
其中$p=0$的项取零；$p>0$而$q=0$的项产生$+\infty$。

令$f=dP/dQ$，则积分等于$\int f\log f\,dQ$。
由不等式
\[
 t\log t\geq-\frac1e\qquad(t\geq0)
\]
与$Q(X)=1$，被积函数的负部积分有下界，不会出现负无穷；正部仍可为$+\infty$，因此定义合法。
再由$t\log t\geq t-1$与$\int f\,dQ=1$，
\[
 D_{\mathrm{KL}}(P\|Q)=\int f\log f\,dQ\geq\int(f-1)\,dQ=0.
\]
等号要求$f=1$在$Q$几乎处处成立，即$P=Q$。
积分合法与积分有限是两件事：绝对连续不保证KL有限。

\paragraph{方向差别与无穷值的例子}

KL的方向不能互换。
对两点概率$P=(1/2,1/2)$、$Q=(3/4,1/4)$，直接代入有限和：
\[
 D_{\mathrm{KL}}(P\|Q)=\frac12\log\frac{2}{3}+\frac12\log 2
  =\frac12\log\frac43,
\]
而
\[
 D_{\mathrm{KL}}(Q\|P)
  =\frac34\log\frac32+\frac14\log\frac12,
\]
两值不同；改变方向同时改变了加权概率和对数比。
KL不对称，也不满足度量的三角不等式，不能把它当作距离使用。

即使互相绝对连续，KL也可能无穷。
在正整数上取
\[
 p_n=\frac{c}{n^2},\quad c=\Bigl(\sum_{n\geq1}n^{-2}\Bigr)^{-1},
  \qquad q_n=(e-1)e^{-n}.
\]
两者都归一化且每个点权重为正，因而互相绝对连续；但
\[
 \log\frac{p_n}{q_n}=n-2\log n+O(1),
\]
充分大的$n$上$p_n\log(p_n/q_n)$至少是某个正常数乘$1/n$，其和发散，故$D_{\mathrm{KL}}(P\|Q)=+\infty$。
这个例子把“零集关系允许定义密度”与“对数比具有足够积分控制”明确分开。

\subsubsection{全变差、Hellinger与KL的关系}
\label{sec:measure-metric-relations}

三种指标建立后，固定同一对概率、同一共同参考，判断它们能相互控制到什么程度。

\paragraph{全变差与Hellinger的双向控制}

对非负数$a,b$，有$(\sqrt a-\sqrt b)^2\leq|a-b|$，因为$|a-b|=|\sqrt a-\sqrt b|(\sqrt a+\sqrt b)$且$\sqrt a+\sqrt b\geq|\sqrt a-\sqrt b|$。
积分得到
\[
 H^2(P,Q)\leq\operatorname{TV}(P,Q).
\]
反向把$|p-q|$拆成$|\sqrt p-\sqrt q|(\sqrt p+\sqrt q)$，用Cauchy–Schwarz：
\[
 \int|p-q|\,d\lambda
 \leq\left(\int(\sqrt p-\sqrt q)^2\,d\lambda\right)^{1/2}
      \left(\int(\sqrt p+\sqrt q)^2\,d\lambda\right)^{1/2}.
\]
第二个因子不超过$2$，而$\int(\sqrt p-\sqrt q)^2\,d\lambda=2H^2$，结合$\operatorname{TV}=\frac12\int|p-q|$得到
\[
 \operatorname{TV}(P,Q)\leq\sqrt2\,H(P,Q).
\]
所以$H$与TV趋零互相等价，但数值和尺度不同。

\paragraph{Pinsker界与反向控制的边界}

自然对数约定下有\term{Pinsker不等式}
\[
 \operatorname{TV}(P,Q)^2\leq\frac12D_{\mathrm{KL}}(P\|Q).
\]
证明思路是先取全变差正差集合$A_+$，把KL按$A_+$及其补集分组；由对数函数的凸性，一般空间上的这一组不小于只在“落入$A_+$/不落入$A_+$”两个事件上的二元KL。
对二元概率，固定第二参数$b$后函数
\[
 a\mapsto a\log\frac ab+(1-a)\log\frac{1-a}{1-b}
\]
的二阶导数为$1/[a(1-a)]\geq4$，且在$a=b$处值与一阶导数为零；积分二阶下界得到该二元KL至少为$2(a-b)^2$，而$|a-b|$正是这两个事件上的全变差，整理即得Pinsker界。
端点情形由连续极限或无穷约定处理，不在边界除以零。

于是KL小先控制全变差，再控制事件概率和有界积分差；Hellinger小也能作同类控制。
反向控制一般不成立：取
\[
 P_\varepsilon=(1-\varepsilon)\delta_0+\varepsilon\delta_1,\qquad Q=\delta_0,
\]
则$\operatorname{TV}(P_\varepsilon,Q)=\varepsilon$，但$P_\varepsilon$在$\{1\}$有正质量而$Q$质量为零，故$D_{\mathrm{KL}}(P_\varepsilon\|Q)=+\infty$。
TV趋零不能推出KL趋零。

\subsection{共同底层距离下的Wasserstein比较}

现在两项概率仍在同一空间，只是比较中增加一个共同距离$d$。
先用联合权重表示两边的质量配对，再用距离幂累计配对成本；不要求预先构造把每个点送到另一个点的映射。

\subsubsection{\term{耦合}与质量配对}

\paragraph{两个边缘与耦合候选}

固定同一度量空间$X$上的Borel概率$P,Q$。
在乘积空间$X\times X$上取联合概率$\gamma$，要求
\[
 \gamma(A\times X)=P(A),\qquad
 \gamma(X\times B)=Q(B)
\]
对所有Borel集合$A,B$成立；满足这两条边缘约束的全体联合概率记为$\Pi(P,Q)$。
一个耦合同时规定从哪里取质量、与哪里的质量配对。
候选集合非空：独立耦合$P\otimes Q$总是可行。
联合权重允许把同一源点的质量分给多个目标点，因此它不是点到点函数；什么时候要求配对由确定性映射实现，会在跨空间一节讨论。

\subsubsection{Wasserstein距离}

\paragraph{距离幂成本与有限矩范围}

取带有给定完备可分度量$d$的空间$X$，本章在这一Polish范围内使用标准结论。
固定$1\leq p<\infty$和基点$x_0$，定义具有有限$p$阶矩的概率集合
\[
 \mathcal P_p(X)=\left\{P:\int_X d(x,x_0)^p\,dP(x)<\infty\right\}.
\]
换基点不改变这个集合：由幂不等式，$d(x,x_0')^p$被$d(x,x_0)^p$的常数倍加一个常数控制，反之亦然。

对$P,Q\in\mathcal P_p(X)$，定义
\[
 W_p(P,Q)
  =\left(\inf_{\gamma\in\Pi(P,Q)}
     \int_{X\times X}d(x,y)^p\,d\gamma(x,y)\right)^{1/p}.
\]
共同底层距离衡量位置差异，耦合约束保证质量守恒，开$p$次方恢复距离尺度；公式中的$x,y$都属于同一空间$X$。
至少存在一个有限成本候选：对独立耦合使用
\[
 d(x,y)^p\leq2^{p-1}\bigl[d(x,x_0)^p+d(y,x_0)^p\bigr],
\]
右侧在$P\otimes Q$下的积分由两项的有限$p$阶矩控制。
候选非空、最优成本有限和最小值由某个耦合达到是三个不同判断，不能从取下确界自动得到后两项。

\paragraph{距离性质与最优计划的范围}

在上述Polish范围与有限矩条件下，$W_p$是$\mathcal P_p(X)$上的距离，并且最优耦合存在。
三角不等式的依据是把具有共同边缘的两份耦合拼接成一份三空间耦合，再用距离的三角不等式和幂估计控制距离幂积分；同一性由$W_p(P,P)=0$与成本为零推出两边缘在对角线上一致。
存在性则要用固定边缘耦合族的紧性以及成本关于\term{弱收敛}的\term{下半连续}性，这里作为标准结果使用，不把简短说明当作完整证明。
最优联合计划不一定由确定性映射给出；若缺少有限$p$阶矩，成本可以为$+\infty$；若没有共同底层距离，也不能直接套用$d(x,y)^p$。

\paragraph{点质量与有限分布的计算}

对$\delta_a,\delta_b$，唯一耦合是$\delta_{(a,b)}$，所以
\[
 W_p(\delta_a,\delta_b)=d(a,b).
\]
这直接显示质量移动的几何距离，而不只是判断两个事件权重是否不同。

对有限离散分布
\[
 P=\sum_i\alpha_i\delta_{x_i},\qquad
 Q=\sum_j\beta_j\delta_{y_j},
\]
任何耦合都由质量矩阵$\gamma_{ij}\geq0$给出，约束为行和、列和守恒：
\[
 \sum_j\gamma_{ij}=\alpha_i,\qquad
 \sum_i\gamma_{ij}=\beta_j.
\]
在这些约束下最小化
\[
 \sum_{i,j}\gamma_{ij}d(x_i,y_j)^p,
\]
再开$p$次方就是$W_p$。
产品矩阵$\alpha_i\beta_j$给出可行点；有限维可行集闭且有界，线性成本连续，所以本例的最小值能达到。
这条有限维论证与一般\term{Polish空间}的存在性结论分别成立。

\subsection{测试函数族与比较方式的对照}

前面分别建立了测度关系、密度差异和运输成本。
本节回到比较的用途：如果最终要读取某些函数的积分，哪一种比较能控制这些读数？
测试函数族给出统一的观察口径；TV和$W_1$在其中有相应表示，但它不是与这些指标互斥的另一类分布。

\subsubsection{指定观测类的最大积分差}

\paragraph{允许的观测与可辨认范围}

固定同一可测空间上的概率$P,Q$和非空实值可测函数类$\mathcal F$，要求每个$f\in\mathcal F$在两侧都绝对可积。
把每个积分看成一种允许读取的观测，定义
\[
 d_{\mathcal F}(P,Q)
  =\sup_{f\in\mathcal F}\left|\int f\,dP-\int f\,dQ\right|,
\]
称为\term{积分概率度量}。
上确界遍历允许的全部观测，不是在源点或耦合上优化。

在共同积分可用的范围内，对称性由绝对差得到，三角不等式逐个测试成立后再取上确界。
结果可以为$+\infty$；只有在所用测度范围内有限、并且函数类能区分不同测度时，它才构成通常意义的距离。
例如在实轴上取
\[
 P=\tfrac12\delta_{-1}+\tfrac12\delta_1,\qquad Q=\delta_0,
\]
测试类只有$f(x)=x$时，两侧一阶读数都为$0$而两项测度不同，小测试类只给出伪距离。

扩大函数类会增加候选观测，上确界不会减小；若允许某个非零积分差的函数任意放大倍数，上确界就成为无穷，所以函数类的尺度限制是定义的一部分。
前面的TV公式正好是两个特例：$0\leq f\leq1$的全部可测函数给出TV，$\|f\|_\infty\leq1$给出两倍TV。
比较数值时不能忽略函数类和归一化。

\subsubsection{Lipschitz观测与一阶Wasserstein}

\paragraph{由耦合成本控制积分差}

固定Polish空间和$P,Q\in\mathcal P_1(X)$，令$g$为$L$-Lipschitz实函数。
它在两侧都可积：选定基点$x_0$后，
\[
 |g(x)|\leq|g(x_0)|+L\,d(x,x_0),
\]
有限一阶矩足以容纳这种不必有界的观测。

对任意$\gamma\in\Pi(P,Q)$，用两个边缘把积分差写成联合积分：
\[
 \int g\,dP-\int g\,dQ
  =\int[g(x)-g(y)]\,d\gamma(x,y).
\]
由Lipschitz条件$|g(x)-g(y)|\leq L\,d(x,y)$，
\[
 \left|\int g\,dP-\int g\,dQ\right|
  \leq L\int d(x,y)\,d\gamma.
\]
这对所有候选耦合成立，对右侧取下确界得到
\[
 \left|\int g\,dP-\int g\,dQ\right|
  \leq L\,W_1(P,Q).
\]
这说明运输比较怎样转化为输出误差控制；此时只证明了一个方向。

\paragraph{Lipschitz测试的对偶表示}

同一范围内还有\term{Kantorovich–Rubinstein对偶}：
\[
 W_1(P,Q)
  =\sup_{\operatorname{Lip}(f)\leq1,\ f(x_0)=0}
     \left|\int f\,dP-\int f\,dQ\right|.
\]
总质量相等允许给测试函数减去任意常数，基点归零不改变积分差，又使函数增长直接由距离控制。
反向等号是标准对偶结论，不是前一段两行估计的推论，这里不伪装成简单证明。

把它放回测试类比较：事件、有界观测和Lipschitz观测保留的信息不同。
后章的核函数空间单位球可以给出另一种测试类，但核均值、特征核条件和MMD构造依赖函数空间理论，留在那里建立；其他密度比散度和KL的变分表示也留给后面的概率分布分析，本章不穷举所有比较公式。

\subsubsection{权重差异与位置差异的对照}

固定实轴Borel空间中的$\delta_a,\delta_b$且$a\neq b$，并列核对：
\[
 \operatorname{TV}(\delta_a,\delta_b)=1,\qquad
 H(\delta_a,\delta_b)=1,\qquad
 D_{\mathrm{KL}}(\delta_a\|\delta_b)=D_{\mathrm{KL}}(\delta_b\|\delta_a)=+\infty,
\]
而
\[
 W_p(\delta_a,\delta_b)=|a-b|.
\]
前几项只使用相同事件或共同参考，最后一项还使用位置间隔；同一底集不意味着所有指标回答同一个问题。
让$b$趋向$a$但始终不重合，几何差异趋于零，固定位置的权重差异却始终为$1$：若关心事件或有界观测误差，应读取TV等控制；若允许小幅位置移动，则读取几何成本，同时核对矩条件。

再取同一Polish空间上的有限加权点概率$\widehat P=\sum_iw_i\delta_{x_i}$和无原子Borel概率$P$。
有限支撑集$S=\{x_i\}$满足$\widehat P(S)=1$、$P(S)=0$，二者相互奇异，故TV恒为$1$；有限离散近似是否有用，不能只要求它在TV下逼近连续目标。
样本或粒子是否在弱测试或几何标准下收敛，还需要后文的抽样和近似条件，本例没有自动给出一致性。
以上比较都把两项测度放在同一可测空间或同一度量空间；下一节允许输入分别来自$X,Y$，先给映射、测度核、共同表示或跨空间成本，再决定实际比较哪两个对象。

\section{不同空间上的映射与比较}
\label{sec:measure-cross-space}

本节研究空间之间的质量怎样传递、联合记录和比较。
确定映射为每个输入指定一个目标点；\term{测度核}为每个输入指定目标空间上的一项非负测度。
二者在相应的可测和质量条件下产生输出测度，再按质量、\term{绝对连续}性或矩条件核对上一节各项比较的适用范围。
联合测度保留输入与输出之间的配对关系：在有限输入测度和核有统一质量界的范围内可以正向构造联合测度；从联合测度反求核则要另查空间条件。
若直接比较原空间之间的配对代价，还要另外给定跨空间成本。
各小节按这些操作各自回答的问题分工，再按定义、运算和比较的依赖安排顺序；输出匹配程度与实现对应所需的代价分别判断。

\subsection{\term{可测映射}的推前与密度变换}

确定映射先回答目标集合接收多少输入质量，再由输出积分公式判断密度怎样变化。
推前不以映射可逆或具有密度为前提；一般定义和复合在本小节集中建立，欧氏可逆换元只回用前面的微积分结论，不重新证明。

\subsubsection{\term{推前测度}与输出积分}

\paragraph{推前的定义与质量边界}

固定可测映射$T:(X,\mathcal A)\to(Y,\mathcal B)$和$X$上的测度$\mu$。
对目标可测集$B\in\mathcal B$定义
\[
 (T_\#\mu)(B)=\mu(T^{-1}(B)).
\]
可测性保证每个目标事件的原像都是$\mathcal A$中的合法集合，从而有权重；这里的$T^{-1}(B)$是集合原像，不要求$T$存在逆函数。

$T_\#\mu$确实是$Y$上的测度。
空集原像为空，质量为零；若$B_n$两两不交，原像$T^{-1}(B_n)$也两两不交且
\[
 T^{-1}\!\left(\bigcup_nB_n\right)=\bigcup_nT^{-1}(B_n),
\]
由$\mu$的可数可加性得到$T_\#\mu$可数可加。
又因为$T^{-1}(Y)=X$，总质量守恒：
\[
 (T_\#\mu)(Y)=\mu(X).
\]
所以\term{有限测度}推前后仍有限，概率推前后仍是概率。

$\sigma$有限性不能无条件保持。
实轴上的勒贝格测度经常值映射$T(x)=0$推前后得到$\delta_0$：目标点$0$的质量为$+\infty$，任何包含它的集合都不能具有有限测度，输出无法由有限质量集覆盖。
若两个带质量源点映到同一目标点，目标点收到的质量是两项权重之和；输出测度由这些集合权重完全确定，它是否在某个参考测度下有密度则要另行判断。

\paragraph{输出观测在输入端的表示}

要计算一般输出观测的积分，先从指标函数做起。
对$g=\mathbf 1_B$，
\[
 \int_Y g\,d(T_\#\mu)
  =(T_\#\mu)(B)=\mu(T^{-1}B)
  =\int_X\mathbf 1_B(Tx)\,d\mu(x)
  =\int_X g(Tx)\,d\mu(x).
\]
沿“指标—非负简单函数—非负可测函数单调逼近”的路线推广，对非负可测$g:Y\to[0,\infty]$有
\begin{equation}
 \int_Y g\,d(T_\#\mu)=\int_X g(Tx)\,d\mu(x),
 \label{eq:measure-pushforward-integral}
\end{equation}
允许两边共同为$+\infty$。
有符号$g$分别处理正负部分，只在绝对可积或能避免无穷相减的合法范围内合并。
无须为一般空间重新证明积分构造。

再取可测映射$S:(Y,\mathcal B)\to(Z,\mathcal C)$。由
\[
 (S\circ T)^{-1}(C)=T^{-1}(S^{-1}(C))
\]
直接得到
\[
 (S\circ T)_\#\mu=S_\#(T_\#\mu):
\]
先把质量送到$Y$、再送到$Z$，与直接送到最终点一致。
若某个可测映射$T:X\to X$满足$T_\#\mu=\mu$，称它\term{保持测度}，这只是输出权重与输入权重相同的特例；本章不展开遍历性或Haar测度存在理论。

\subsubsection{输出密度与质量重组}

\paragraph{回用欧氏换元计算输出密度}

式~\eqref{eq:measure-pushforward-integral}把目标观测复合到输入端，不要求输出具有密度，也不是在原空间乘一个Radon–Nikodym系数。
给定参考体积和映射的正则性后，才谈得上输出密度。

设$U,V\subseteq\mathbb R^d$是同维欧氏开集，$T:U\to V$是$C^1$微分同胚，输入测度$\mu$相对于$U$上的勒贝格体积有非负密度$p$。
这里要求全局一一对应、逆映射光滑；局部可逆不足以替代这些假设。
对非负测试函数$g$，先用推前积分恒等式，再回用多元换元公式：
\[
 \int_V g(y)\,d(T_\#\mu)(y)
  =\int_U g(Tx)p(x)\,dx
  =\int_V g(y)\,p(T^{-1}y)\,|\det DT^{-1}(y)|\,dy.
\]
由于这对所有非负$g$成立，输出密度的几乎处处表示为
\[
 q(y)=p(T^{-1}y)\,|\det DT^{-1}(y)|.
\]
利用逆导数关系$DT^{-1}(Tx)=(DT(x))^{-1}$，也可写成
\[
 q(Tx)\,|\det DT(x)|=p(x).
\]
体积伸缩因子与密度系数相乘正好保持质量；有符号测试还需积分合法性。
这里只衔接一般推前和已有换元结论，不重做微积分中的证明。
平移映射给出$q(y)=p(y-b)$；可逆线性映射$T(x)=Ax$给出
\[
 q(y)=\frac{p(A^{-1}y)}{|\det A|};
\]
负比例缩放由行列式绝对值处理体积。
多原像、降维和退化映射不满足全局微分同胚条件，必须改用原像计算，不能挑选某一个逆分支硬套公式。

\paragraph{合并、投影与原子输出}

离散合并时，目标点收到全部原像源点的权重之和；连续映射若有多个逆分支，也要把所有原像分支的贡献相加。
推前本身始终有定义，但目标密度不能由单个全局逆自动保证。

设$\mathbb R^2$上的非负可积密度为$p(x,y)$（对乘积可测结构可测；若只在勒贝格完备化下可测，先取几乎处处相同的乘积可测代表），考虑投影$T(x,y)=x$。
对一维Borel集$B$，原像为$B\times\mathbb R$，由Tonelli定理，
\[
 (T_\#\mu)(B)
  =\int_{B\times\mathbb R}p(x,y)\,dx\,dy
  =\int_B\left(\int_{\mathbb R}p(x,y)\,dy\right)dx.
\]
因此边缘密度
\[
 q(x)=\int_{\mathbb R}p(x,y)\,dy
\]
正是该投影的输出密度，目标参考是一维勒贝格测度。
这是已有累次积分的用途，不再铺开构造。

改为常值映射$T(x,y)=y_0$时，有限输入总质量$m$被推为原子$m\delta_{y_0}$，因为目标集合的原像只能是全空间或空集。
在正维欧氏目标空间中且$m>0$时，它相对于勒贝格测度没有普通密度；$m=0$则给出零测度，要与非零原子分开。
投影和常值输出说明：不可逆映射可以保持密度，也可以产生原子，两种情形都可能合并原本不同的点。
确定映射仍要求每个输入只对应一个输出；若同一输入还要把质量分给多个目标，就要由下一小节的测度值规则取代单个输出点。

\subsection{测度核与\term{测度值映射}}

本小节把“每个输入对应一个目标点”扩展为“每个输入对应目标集合上的一项非负测度”。
新对象仍是映射，但输出从点换成测度；随输入变化的集合质量必须可测，才能对输入测度累计。
先给一般测度核的定义，再为积分、复合和联合构造明确采用有限输入测度及统一有限质量界的核；这是本章选择的充分工作范围，不作为这些运算的最弱条件。
总质量为$1$是额外的归一化条件，不写入一般核的定义。

\subsubsection{随输入变化的输出权重}

\paragraph{测度核的两项条件}

先看有限输入、输出集合上的非负矩阵：固定一行，矩阵元素给出该输入对应各目标点的质量，对一组目标点求和便得到集合质量。
不要求各行之和相同或等于$1$；当输入扩展到一般集合时，还必须让每个目标集合的质量随输入可测，才能对输入测度累计。

\begin{definition}[测度核]
\label{def:measure-kernel}
设$(X,\mathcal A)$与$(Y,\mathcal B)$为可测空间。映射
\[
 K:X\times\mathcal B\to[0,\infty]
\]
若满足：
\begin{enumerate}
 \item 对每个$x\in X$，$K_x=K(x,\cdot):\mathcal B\to[0,\infty]$是$(Y,\mathcal B)$上的非负测度；
 \item 对每个$B\in\mathcal B$，函数$x\mapsto K(x,B)$关于$\mathcal A$可测，
\end{enumerate}
就称$K$为从$X$到$Y$的测度核，记作$K:X\leadsto Y$。
\end{definition}

第一项约束目标端的集合赋权，第二项约束输入端的可测依赖，缺一不可。
这里的测度核与函数空间章中表示相似度的正定核是完全不同的对象。

一般核定义本身没有质量上界。
本章主要运算进一步要求存在与$x$无关的常数$M<\infty$，使
\[
 K(x,Y)\leq M\qquad\text{对所有 }x\in X,
\]
称$K$具有统一有限质量界。
逐\term{点质量}$K(x,Y)<\infty$并不保证存在共同上界，两者不能混用。
若还满足
\[
 K(x,Y)=1\qquad\text{对所有 }x,
\]
称$K$为\term{归一化测度核}（马尔可夫核）；这一附加条件用于保持总质量和边缘。

对可测映射$T:X\to Y$，取
\[
 K_T(x,B)=\mathbf 1_B(Tx)=\delta_{T(x)}(B).
\]
每个$x$对应点质量，每个$B$对应可测的示性函数，所以它是归一化测度核。
确定推前因此成为测度核传递的特例；一般核允许一项输入测度在多个目标点或集合之间分配质量，其总质量还可以随$x$变化。

\subsubsection{核积分与核复合}

\paragraph{输出测度与观测函数的传递}

给定$X$上的有限非负测度$\mu$和满足统一质量界$K(x,Y)\leq M<\infty$的核$K$，定义输出测度
\[
 (\mu K)(B)=\int_X K(x,B)\,\mu(dx).
\]
对两两不交的$B_n$，用核在固定$x$处的可数可加性和非负单调收敛，可把求和移出积分：
\[
 (\mu K)\!\left(\bigcup_nB_n\right)
  =\int_X\sum_nK(x,B_n)\,\mu(dx)
  =\sum_n\int_XK(x,B_n)\,\mu(dx),
\]
故$\mu K$是测度。
总质量满足
\[
 (\mu K)(Y)=\int_XK(x,Y)\,d\mu\leq M\mu(X).
\]
归一化核保证$(\mu K)(Y)=\mu(X)$，一般核不保证质量守恒。

对非负可测$g:Y\to[0,\infty]$，定义
\[
 Kg(x)=\int_Y g(y)\,K(x,dy).
\]
$g$为指标时，$Kg$就是核定义中的可测函数；简单函数对应有限线性组合；非负单调逼近保证$Kg$可测，并得到核积分恒等式
\begin{equation}
 \int_Y g\,d(\mu K)=\int_X Kg\,d\mu,
 \label{eq:measure-kernel-integral}
\end{equation}
允许两边共同为$+\infty$。
这里必须用到核随输入可测这第二项条件，不能只说每个截面的积分存在。
对有界可测实函数$g$，
\[
 |Kg(x)|\leq M\|g\|_\infty,
\]
各点积分有限。
一般有符号$g$在$\int|g|\,d(\mu K)<\infty$时，先对$|g|$使用式~\eqref{eq:measure-kernel-integral}，得到
\[
 \int_Y|g(y)|\,K(x,dy)<\infty
\]
在$\mu$几乎处处成立；在剩余零集上为$Kg$选定可测值，再合并正负部分即可，避免无穷相减。

用一个离散例复算。
输入点为$a,b$，$\mu(\{a\})=1$、$\mu(\{b\})=2$；输出点为$0,1$，核两行分别为$(2,1)$和$(0,2)$。
输出在$0,1$处的质量分别为
\[
 1\cdot2+2\cdot0=2,\qquad 1\cdot1+2\cdot2=5,
\]
总质量为$7$，等于按输入权重累计的行质量$1\cdot3+2\cdot2=7$，但不等于输入总质量$3$。
取$g(y)=y$，两行核积分分别为$1$和$2$，输入积分为$1\cdot1+2\cdot2=5$，与输出积分一致。

\paragraph{连续两次传递的复合}

给定核$K:X\leadsto Y$与$L:Y\leadsto Z$，统一质量界分别为$M,N<\infty$。
前一个核为每个$x$给出$Y$上的测度，后一个核以$Y$中的点为输入，按前一测度累计后一核的输出，定义
\[
 (KL)(x,C)=\int_Y L(y,C)\,K(x,dy).
\]
对固定$C$，被积函数关于$y$可测；对不交的$C_n$使用非负单调收敛得到可数可加。
质量界为
\[
 (KL)(x,Z)\leq N\,K(x,Y)\leq NM,
\]
所以复合核仍在统一有限质量界范围内；两核都归一化时复合也归一化，一般情况下不能把总质量写成$1$。

先对目标指标使用核积分恒等式，得到测度恒等式
\[
 (\mu K)L=\mu(KL);
\]
再对简单及非负函数推广，得到逐点的扩展积分恒等式
\[
 (KL)g=K(Lg).
\]
有符号$g$在$\int|g|\,d[\mu(KL)]<\infty$时，按前述方式选取可测积分版本，观测恒等式在$\mu$几乎处处意义下成立。
两个式子分别说明：先后传递输出测度与先后传回观测是一致的操作。
把$K,L$依次取为可测映射$T,S$的点核时，
\[
 K_TK_S=K_{S\circ T},
\]
恢复推前的复合式，说明按积分定义的核复合扩展了点值映射的普通复合。
本章只建立有限串联的通用运算，不展开无限串联或路径空间理论。

\subsubsection{测度值映射的可测性}

\paragraph{用集合质量检验测度族的可测依赖}

记$(Y,\mathcal B)$上全部非负测度组成的集合为$\mathcal M_+(Y)$，在其上取由评价映射
\[
 e_B(\nu)=\nu(B),\qquad B\in\mathcal B
\]
生成的$\sigma$代数，评价值取$[0,\infty]$。
在这个结构下，映射$F:X\to\mathcal M_+(Y)$可测，当且仅当对每个$B\in\mathcal B$，函数
\[
 x\longmapsto F(x)(B)
\]
都可测。
测度核的两项条件正好使$F(x)=K_x$成为这样一项测度值可测映射；统一有限质量界和归一化只进一步限制值域，不改变可测性判断。

以$Y^n$的乘积可测结构为输入，固定非负有限权重$a_1,\ldots,a_n$，定义
\[
 F(y_1,\ldots,y_n)=\sum_{i=1}^na_i\delta_{y_i}.
\]
对任意$B$，评价为$\sum_i a_i\mathbf 1_B(y_i)$，是有限个坐标可测函数之和，因此加权点测度关于输入可测，总质量有共同界$\sum_i a_i$。
取$a_i=1/n$得到\term{经验测度}这一特例。
这里要区分两个层面：每个输入给出一项测度只检查了目标端；输入到测度的映射可测还要检查这些评价函数。
当前直接指定评价$\sigma$代数，不把它无条件等同于任意空间上弱拓扑的Borel结构。
完成从点值映射到测度值映射的可测接口后，就可以形成联合测度，再判断反向分解是否存在。

\subsection{联合测度与测度\term{核分解}}

输出测度只保留目标端质量；若还要记录哪部分输入贡献了哪些输出，应保留联合测度。
本小节分别回答：给定输入测度和核时怎样构造联合权重；给定联合权重时何时能反求核。
正向构造沿用有限输入与统一质量界，反向分解则从第一边缘出发，需要额外的空间条件。
\term{条件积分}只作为理解反向分解的通用基础，其进一步代数法则另行展开。

\subsubsection{由输入测度和核构造联合测度}

\paragraph{联合权重、两个边缘与乘积特例}

给定$X$上的有限非负测度$\mu$和具有统一有限质量界的核$K:X\leadsto Y$。
先在可测矩形上规定
\[
 J(A\times B)=\int_A K(x,B)\,\mu(dx),
\]
再对乘积可测集合$E\subseteq X\times Y$用
\[
 J(E)=\int_X\left[\int_Y\mathbf 1_E(x,y)\,K(x,dy)\right]\mu(dx)
\]
构造联合测度。
由矩形指标进入，利用核的统一有限质量界把有界函数的结论经单调类推广，得到截面积分关于$x$可测；再以单调收敛核对可数可加。
总质量受$M\mu(X)$控制。
这个构造没有假定$J$在某个乘积参考测度下有密度。

在矩形上分别取$B=Y$与$A=X$，得到两个边缘：
\[
 (\pi_X)_\#J(A)=\int_A K(x,Y)\,d\mu(x),
 \qquad
 (\pi_Y)_\#J=(\mu K).
\]
第一边缘是在$\mu$上乘以行总质量$K(x,Y)$得到的测度；核归一化时它等于原输入$\mu$，一般情形不保证。
对非负乘积可测$h$，沿指标、简单函数和单调逼近有
\begin{equation}
 \int_{X\times Y}h\,dJ
  =\int_X\int_Y h(x,y)\,K(x,dy)\,\mu(dx);
 \label{eq:measure-joint-integral}
\end{equation}
有符号情形要求$\int|h|\,dJ<\infty$。
记号$\mu(dx)K(x,dy)$表示这条联合积分规则。

两个特例值得写出。
若$K(x,\cdot)=\nu$是与$x$无关的固定有限测度，矩形质量变为$\mu(A)\nu(B)$，由矩形生成族上的唯一性得$J=\mu\otimes\nu$，第一边缘为$\nu(Y)\mu$。
若$K=K_T$是可测映射$T$的归一化点核，联合测度是\term{图测度}
\[
 J=(\operatorname{id},T)_\#\mu,
\]
第一边缘为$\mu$。
乘积测度和确定图测度都是联合构造的特例，但其边缘要按各自总质量核对；反向寻找描述给定联合关系的核，需要下面的测度分解。

\subsubsection{从联合测度反求测度核}

\paragraph{子$\sigma$代数上的条件积分与观测值表示}

反向分解要先理解“保留哪些输入信息”如何表示积分。
固定有限测度空间$(\Omega,\mathcal F,\mu)$、子$\sigma$代数$\mathcal G\subseteq\mathcal F$和可积实函数$h$。
只在$\mathcal G$的集合上考察
\[
 \nu(A)=\int_A h\,d\mu,\qquad A\in\mathcal G,
\]
得到相对于\term{限制测度}$\mu|_{\mathcal G}$绝对连续的有限\term{有符号测度}；限制可测集合族保持底集不变，不能说成把$\Omega$缩小。

用$h^+$和$h^-$定义两个受限的有限非负测度
\[
 \nu_\pm(A)=\int_A h^\pm\,d\mu,
\]
它们都对$\mu|_{\mathcal G}$绝对连续。
对两者分别应用非负版本的Radon–Nikodym定理，取得$\mathcal G$可测、可积的密度$H_\pm$，令$H=H_+-H_-$的有限可测版本，则
\[
 \int_A H\,d\mu=\int_A h\,d\mu\qquad(A\in\mathcal G),
\]
并且$H$在$\mu$几乎处处意义下唯一。
这个条件积分保存每个可观测集合上的累计，不要求$H=h$逐点成立；有限测度使受限参考仍有限，不无条件推广到任意$\sigma$有限空间的子$\sigma$代数。

若$\mathcal G=\sigma(S)$由可测观测$S:\Omega\to X$生成，实值$\mathcal G$可测函数可以写成$H=\phi\circ S$，其中$\phi:X\to\mathbb R$可测。
理由是：简单的$\mathcal G$可测函数的层集合是$S$的原像，可写成$\phi(S)$；一般实值可测函数由简单函数逐级逼近，在目标端不收敛处选定一个可测值即可。
因此“依赖观测的信息”能表示成观测值的函数，而不是在每条纤维$S^{-1}(x)$上任意取值。

把这两步用于非零有限联合测度$J$：取$\Omega=X\times Y$、$\mu=J$、$S=\pi_X$、$\mathcal G=\sigma(\pi_X)$，对每个目标集合$B$令$h(x,y)=\mathbf 1_B(y)$。
条件积分因子化为$H_B=k_B\circ\pi_X$；用第一边缘$\mu_X=(\pi_X)_\#J$的输出积分公式，得到
\[
 \int_A k_B\,d\mu_X=J(A\times B)\qquad(A\in\mathcal A).
\]
右端介于$0$和$\mu_X(A)$之间，所以可选取$[0,1]$值的可测代表$k_B$，它说明输入端的条件质量函数怎样取得。
但逐个集合取得$k_B$，还没有保证这些代表在同一组输入上组成可数可加、总质量为$1$的测度；把一族条件积分协调成归一化测度核，需要下一目在标准Borel范围内的存在定理。

\paragraph{标准Borel范围内的联合分解}

反向分解当前采用两侧都是\term{标准Borel空间}的充分范围，不宣称这是最弱条件。
标准Borel结构提供可数生成族、决定测度以及协调各版本所需的条件，它没有给定跨空间距离。

\begin{theorem}[测度的核分解]
\label{thm:measure-disintegration}
设$X,Y$为标准Borel空间，$J$是$X\times Y$上的非零有限测度，第一边缘为
\[
 \mu_X=(\pi_X)_\#J.
\]
则存在归一化测度核$K:X\leadsto Y$，使
\[
 J(A\times B)=\int_A K(x,B)\,\mu_X(dx),
\]
从而
\[
 J=\mu_X(dx)\,K(x,dy)
\]
且第二边缘为$\mu_XK$。
\end{theorem}

归一化是“按实际第一边缘分解”带来的，不要求$J$的总质量为$1$；也不能反推任意给定输入和核时第一边缘必然不变。
存在性证明需要把可数决定族上的条件积分协调成同一个测度版本，本章陈述准确结论并说明这一构造障碍，完整证明另行展开。

核的唯一性是共同$\mu_X$零集以外的测度相同：若$K,L$都给出该分解，则存在同一个可测集合$N$满足$\mu_X(N)=0$，使对每个$x\notin N$和全部$B\in\mathcal B$都有
\[
 K(x,B)=L(x,B).
\]
标准Borel的可数决定族先允许合并可数个零例外，再由有限测度的唯一性推广到全部可测集合；不能把“逐集合几乎处处相等”直接写成无条件的共同零集结论。

取$X=Y=[0,1]$的Borel空间、输入测度为区间上的勒贝格测度$\lambda$、$T(x)=x$，联合图测度
\[
 J=(\operatorname{id},T)_\#\lambda
\]
集中在对角线上。
它的两个边缘都是$\lambda$，但$(\lambda\otimes\lambda)$给对角线赋零而$J$给对角线质量$1$，所以$J$相对于乘积参考没有密度；核$K(x,\cdot)=\delta_x$仍逐项重建$J$。
这个例子区分“没有联合密度”和“没有测度核分解”，也说明坐标分解与上一节相对于参考测度的\term{Lebesgue分解}回答的是不同问题。

\subsection{建立共同空间后的测度比较}

本小节先确定两项输出在哪里比较，再判断比较结果保留了哪些输入区别。
给定目标映射或共同表示，都是让同空间指标具有合法输入的方式；对输入差异的控制则以同一输出规则为前提。
联合和\term{耦合}构造用来说明可见输出与原输入的关系，不能把表示相同误写成原对象相同。

\subsubsection{目标空间与共同表示}

\paragraph{指定映射后的目标端差异}

给定$X$上的概率$P$、$Y$上的概率$Q$和预先确定的可测映射$T:X\to Y$，先形成输出测度
\[
 R=T_\#P,
\]
再在$Y$上比较$R$与$Q$。
同一个目标事件$B$对应源事件$T^{-1}(B)$，所以比较的是映射后输出对目标事件的质量判断，不能直接对原来不同空间的$P,Q$写密度相减。
若给定的是归一化测度核，则以$R=PK$形成相同的目标端问题；一般核可能改变总质量，必须先按所选比较方式核对输出范围。

若采用TV或Hellinger，在$Y$上取共同有限参考$\lambda_Y=R+Q$即可得到两项\term{Radon–Nikodym密度}，不必预设勒贝格密度。
若采用$D_{\mathrm{KL}}(R\|Q)$，要固定方向，检查$R\ll Q$和对数比的可积性；绝对连续不足以保证KL有限。
零集关系、Lebesgue分解和测试函数比较都在同一个目标可测空间上判定。

若采用Wasserstein比较，要在Polish目标空间$(Y,d_Y)$的Borel结构上取$R,Q$，固定$1\leq p<\infty$，并要求$R,Q\in\mathcal P_p(Y)$。
确定映射时，$R$的矩由输出积分写为
\[
 \int_Y d_Y(y,y_0)^p\,dR(y)
  =\int_X d_Y(Tx,y_0)^p\,dP(x);
\]
归一化测度核时则写为
\[
 \int_X\int_Y d_Y(y,y_0)^p\,K(x,dy)\,P(dx).
\]
源空间在另一个距离下的有限矩不能无条件替代目标矩。
例如以$[0,1]$上的均匀概率为输入，令
\[
 T(0)=0,\qquad T(x)=\frac1x\quad(x>0),
\]
输入在原距离下具有所有有限阶矩，但对任意$1\leq p<\infty$，
\[
 \int_0^1 T(x)^p\,dx=\int_0^1x^{-p}\,dx=+\infty,
\]
输出不属于Wasserstein有限距离的范围。

\paragraph{在共同表示空间中比较}

两侧都要重新表示时，固定可测映射$S:X\to Z$和$T:Y\to Z$，比较$S_\#P$与$T_\#Q$。
先要说明$Z$的可测结构以及两个映射各保留了什么输入区别，随后所有事件、密度和测试函数都取在这个共同空间内。
测度值规则也可由两侧各自到$Z$的核形成输出；沿用当前的归一化指标时，须核对两核都归一化；使用Wasserstein时还要检查输出矩条件，不能由一般核的可测性自动推出。

对权重差异，在$Z$上取共同参考$S_\#P+T_\#Q$，回用TV、Hellinger或指定方向的KL。
对几何差异，取$Z$为Polish度量空间$(Z,d_Z)$，要求两个表示对其Borel结构可测，固定$1\leq p<\infty$，检查两侧推前都有有限$p$阶距离矩，再计算
\[
 W_p^Z(S_\#P,T_\#Q).
\]
$X,Y$上各自的距离$d_X,d_Y$有定义，并没有给出跨空间的$d(x,y)$。
取$Z=Y$、$S$为源到目标的映射、另一侧取恒等映射，就恢复直接的目标端比较。
一般共同表示的数值依赖$Z,S,T$和所选指标；若把两侧表示都取成同一个常值，所有输入都落到同一目标点，零差异不能说明原对象相同。

\subsubsection{输出规则与可辨认范围}

\paragraph{确定映射下的信息丢失}

在同一源空间$\mathbb R$上取$P_1=\delta_{-1}$、$P_2=\delta_1$，用$T(x)=|x|$映到$[0,\infty)$。
两个源测度在事件$\{-1\}$上的质量不同，但推前后都是$\delta_1$，所以目标端的TV、Hellinger和可用的Wasserstein差异都为零；输出相同只说明这个表示看不到符号区别。

对同一可测$T$和同一源空间上的任意概率$P_1,P_2$，逐个目标事件有
\[
 |(T_\#P_1)(B)-(T_\#P_2)(B)|
  =|P_1(T^{-1}B)-P_2(T^{-1}B)|
  \leq\operatorname{TV}(P_1,P_2).
\]
对$B$取上确界得到TV在可测推前下不增：
\[
 \operatorname{TV}(T_\#P_1,T_\#P_2)\leq\operatorname{TV}(P_1,P_2),
\]
因为只检测由目标原像组成的那部分源事件。

两侧使用不同表示时不能直接套用这个估计。
即使源测度都是$\delta_0$，取$S(0)=0$、$T(0)=1$也得到$\delta_0$与$\delta_1$，TV为$1$。
同一映射的非扩张性不能改写成任意两个表示的非扩张性；粗观测也不能反向控制全部输入事件。

\paragraph{同一归一化测度核下的TV控制}

固定同一源空间上的概率$P,Q$和同一个到$Y$的归一化测度核$K$，于是$PK,QK$仍是概率。
对每个目标事件$B$，输出差为
\[
 |(PK)(B)-(QK)(B)|
  =\left|\int_XK(x,B)\,d(P-Q)(x)\right|.
\]
函数$x\mapsto K(x,B)$可测，且由核的非负性和归一化有$0\leq K(x,B)\leq1$。
回用TV关于$[0,1]$值测试函数的精确公式，逐事件得到该差不超过$\operatorname{TV}(P,Q)$，再对$B$取上确界：
\[
 \operatorname{TV}(PK,QK)\leq\operatorname{TV}(P,Q).
\]
若只调用$\|f\|_\infty\leq1$的一般有界函数界会多出系数$2$；这里正好可用核函数的非负和单位上界取得精确常数。
点核特例恢复确定映射的估计。
一般核可能改变总质量，以上归一化结论不直接适用；把一个归一化核乘常数$2$会使输出总质量加倍，已不在概率TV范围内。
两侧改用不同核时也要另行控制输出规则差异，前面的两个不同常值映射已经给出反例。
本处只证明同核TV控制，不宣布所有散度和几何指标在任意核下都收缩；KL的数据处理解释留给概率分布分析。

\paragraph{共同输入构造耦合上界}

为把输出表示连到运输比较，可取共同可测空间$(\Omega,\mathcal F)$上的概率$R$和可测映射$F:\Omega\to X$、$G:\Omega\to Y$。
$(F,G)$对乘积$\sigma$代数可测，联合推前
\[
 \gamma=(F,G)_\#R
\]
的两个边缘分别是$F_\#R$和$G_\#R$，所以它给出两项输出的一个耦合。

若两项输出都在同一个Polish度量空间$(Z,d)$中，固定$1\leq p<\infty$且两边都有有限$p$阶矩，把这个可行耦合代入Wasserstein定义：
\[
 W_p^p(F_\#R,G_\#R)
  \leq\int_\Omega d(F(\omega),G(\omega))^p\,R(d\omega).
\]
中间由推前积分计算该耦合的成本，再由下确界不超过任一候选成本得到上界。
共同输入在应用中可以是共享噪声，但这里只需可测映射和给定的$R$；该构造给出可行候选和成本上界，不保证最优，也不表明任意给定的两边缘都自带这种表示。

\subsection{给定跨空间成本的运输比较}

本小节保持源和目标各自的空间，直接规定一对源点、目标点配对要付出什么成本。
共同空间比较给出输出端差异；跨空间运输把质量守恒作为约束，在所有允许对应中累计并优化给定成本。
两者不能只因为都出现映射或耦合就等同。

\subsubsection{乘积空间上的耦合与成本}

\paragraph{跨空间质量匹配与成本累计}

把$X\times X$上的耦合推广到$X\times Y$：给定$X,Y$上的概率$P,Q$，以$\Pi(P,Q)$表示第一边缘为$P$、第二边缘为$Q$的乘积空间概率$\gamma$，即
\[
 \gamma(A\times Y)=P(A),\qquad
 \gamma(X\times B)=Q(B).
\]
两个约束分别固定出发质量与到达质量，不要求两个底层空间相同。

另行给定对乘积结构可测的非负成本函数
\[
 c:X\times Y\to[0,\infty],
\]
定义最优运输成本
\[
 \mathcal C_c(P,Q)
  =\inf_{\gamma\in\Pi(P,Q)}
    \int_{X\times Y}c(x,y)\,d\gamma(x,y).
\]
非负累计允许取$+\infty$；$c$规定的是跨空间配对代价，并不是自动存在的一条共同距离。

当两侧在同一度量空间、$c(x,y)=d(x,y)^p$，并满足前面$1\leq p<\infty$和有限$p$阶矩条件时，这个最优成本就是$W_p^p$。
一般跨空间成本的最优值不自动满足对称性、同一性或三角不等式，也不能只开$p$次方就称为\term{Wasserstein距离}。
即使两侧恰为同一空间，取$c\equiv0$也会使不同概率之间的最优成本一律为零，已不能满足距离的分离要求。

\subsubsection{运输映射与\term{运输计划}}

\paragraph{不可拆分的映射与可拆分的计划}

若要求每个源点只送到一个目标点，候选是满足
\[
 T_\#P=Q
\]
的可测映射$T:X\to Y$，其成本为
\[
 \int_X c(x,T(x))\,dP(x),
\]
这给出\term{Monge问题}。
\term{Kantorovich问题}使用全部耦合，允许把同一源点的质量拆给不同目标；两类候选遵守同样的边缘守恒要求。
标准Borel范围内的计划还可由联合分解写成
\[
 P(dx)\,K(x,dy),
\]
这里的$K$由实际第一边缘$P$分解得到，故归一化，输出约束成为$PK=Q$；点核对应确定映射。
不能把任意未归一化的核直接当作满足双边缘守恒的运输计划。

取离散源$X=\{0\}$、目标$Y=\{-1,1\}$，幂集可测，令
\[
 P=\delta_0,\qquad Q=\tfrac12\delta_{-1}+\tfrac12\delta_1.
\]
确定映射只能把全部源质量送到单个$T(0)$，没有可行运输映射；计划
\[
 \gamma=\tfrac12\delta_{(0,-1)}+\tfrac12\delta_{(0,1)}
\]
则满足两个边缘约束。
若$c(0,-1)=c(0,1)=1$，该计划成本为$1$；其核在唯一源点把质量各分一半，具体说明计划放宽了什么限制。
对任何可行映射$T$，图计划$\gamma_T=(\operatorname{id},T)_\#P$具有所需边缘，推前积分给出$\int c\,d\gamma_T=\int c(x,T(x))\,dP$，所以映射候选嵌入计划候选；并非每个计划都来自确定映射，最优计划存在也不替代最优映射存在。

\subsubsection{成本判定与比较范围}

\paragraph{可行、有限与达到分别判断}

产品耦合$P\otimes Q$始终是一个耦合，所以Kantorovich计划的可行域非空；这与原子分裂例中Monge映射可行域为空形成对照。
非空只回答是否存在守恒配对，不回答配对成本是多少。

要得到有限最优成本，至少要有一个有限成本候选。
例如
\[
 \int c\,d(P\otimes Q)<\infty
\]
足以提供这样的候选，但不是唯一途径；若$c$处处为$+\infty$，所有耦合的成本都是无穷。
有限成本条件与概率总质量为$1$是两件事，后者不能控制任意给定$c$的累计。

最小值是否达到还要另查紧性和成本正则性。
对两个\term{Polish空间}上的Borel概率，在非负\term{下半连续}成本及存在有限成本候选等标准条件下，可调用最优计划存在结论；一般可测成本不能无条件套用。
紧性怎样提供\term{弱收敛}子列、下半连续成本怎样给出积分下界，将在\ref{sec:measure-tightness}一目补齐；这里不由计划可行性推出映射存在。

\paragraph{目标匹配、运输代价与结构比较边界}

用单点源和单点目标区分两类比较。
取$X=\{a\}$、$Y=\{b\}$、$P=\delta_a$、$Q=\delta_b$、$T(a)=b$，并令$c(a,b)=7$。
目标端$T_\#P=Q$，同空间差异为零，但唯一跨空间配对的运输成本为$7$。
输出匹配问的是是否得到所需权重，运输成本问的是实现这种对应当价多大；前者为零不意味着后者很小。

因此比较前必须先给定映射、核、共同表示或跨空间成本；仅有$(X,d_X)$和$(Y,d_Y)$并没有规定跨空间点对代价$c(x,y)$。
本章按这些已经给定的对应进行测度比较，还没有从两侧内部距离自行构造空间之间的内在结构距离；这种结构比较的进一步理论另行建立。

\section{测度的收敛与极限操作}
\label{sec:measure-convergence}

本节把前面的比较量用于测度序列和时间参数曲线。
先区分收敛方式，再建立弱极限下事件、映射和成本的可用判据，随后说明什么时候存在极限子列，最后简述时间参数曲线。
各项分别承担定义、稳定操作、极限存在和应用的职责，不把它们混作一种收敛分类。
比较始终发生在固定空间中；变空间对象先按上一节给出共同表示。

\subsection{测度序列的收敛方式}
\label{sec:measure-sequence-convergence}

同一概率列可以由不同观测判定是否趋近：全部事件的统一误差、有界连续观测量，或者加入距离矩的尾部控制。
每种方式的对象、量词和所需附加结构都要分别写清。

\subsubsection{\term{逐集收敛}与全变差收敛}

\paragraph{从比较误差定义趋近}

固定概率列$P_n$和候选概率$P$。
最强的常用事件判据是
\[
 \operatorname{TV}(P_n,P)\to0:
\]
所有可测事件的误差用同一个趋于零的量统一控制。
由\ref{sec:measure-metric-relations}小节的双向不等式，这等价于$H(P_n,P)\to0$；$D_{\mathrm{KL}}(P_n\|P)\to0$是取得这种控制的充分保证（Pinsker界），但由$P_\varepsilon=(1-\varepsilon)\delta_0+\varepsilon\delta_1$与$\delta_0$的例子可知不能反推。

更弱的事件判据是逐集收敛：对每个固定事件$A\in\mathcal A$，
\[
 P_n(A)\to P(A).
\]
给定误差精度后，起始指标$N$允许依赖事件$A$；全变差收敛则要求同一个$N$对全部事件有效。
所以全变差收敛推出逐集收敛，反之一般不成立。
对固定的有界可测函数$g$，逐集收敛仍能控制其积分：先用有限简单函数在$\|g\|_\infty$意义下一致逼近$g$，再用概率总质量为$1$控制两端逼近误差，得到
\[
 \int g\,dP_n\to\int g\,dP.
\]

\subsubsection{\term{弱收敛}与连续测试}

\paragraph{弱收敛与位置移动}

在固定度量空间的Borel概率上，定义$P_n$弱收敛到$P$，记为$P_n\Rightarrow P$，若对每个有界连续函数$\varphi$都有
\[
 \int\varphi\,dP_n\to\int\varphi\,dP.
\]
测试类由“全部可测事件”缩小为“有界连续观测量”。
由于事件指标一般不连续，逐集收敛能推出弱收敛（对每个固定有界连续函数使用上面的积分结论），但弱收敛不自动控制不连续指标，也不自动控制无界矩。

在实轴上取
\[
 P_n=\delta_{1/n},\qquad P=\delta_0.
\]
对任意有界连续$\varphi$，测试积分就是$\varphi(1/n)\to\varphi(0)$，故$P_n\Rightarrow P$；但固定事件$\{0\}$的质量差始终为$1$，\term{全变差距离}也始终为$1$。
同一个位置移动的例子说明，采用哪种比较标准会改变“趋近”的结论。

\subsubsection{Wasserstein收敛与矩条件}

\paragraph{弱趋近与距离尾部的共同控制}

在本章的\term{Polish空间}范围内，$1\leq p<\infty$，$P_n,P\in\mathcal P_p(X)$，Wasserstein收敛有标准判据：
\begin{equation}
 W_p(P_n,P)\to0
 \quad\Longleftrightarrow\quad
 P_n\Rightarrow P
 \ \text{且}\
 \int d(x,x_0)^p\,dP_n\to\int d(x,x_0)^p\,dP.
 \label{eq:measure-wasserstein-criterion}
\end{equation}
这里把判据作为标准结论使用，下面解释矩条件限制的究竟是什么。

令$z(x)=d(x,x_0)^p$。
弱收敛和全矩收敛成立时，对有界连续截断$\min(z,R)$先用弱收敛：
\[
 \int\min(z,R)\,dP_n\to\int\min(z,R)\,dP.
\]
再从全矩中减去有限截断部分，得到
\[
 \int(z-R)_+\,dP_n\to\int(z-R)_+\,dP.
\]
先取足够大的$R$使极限尾部小，再让$n$足够大，并用有限个首项各自可积处理剩余项，就能取得对全列的共同尾部控制。
利用
\[
 z\,\mathbf 1_{\{z>2R\}}\leq2(z-R)_+,
\]
共同超额控制转化为真正的尾积分控制。
反向则由截断积分收敛和共同尾部控制恢复全矩收敛。
所以弱收敛之外还需要距离尾部的一致可积性，而不只是每一项矩分别有限。

\paragraph{小概率远距离的反例}

在实轴上取$n\geq2$及
\[
 P_n=(1-n^{-p})\delta_0+n^{-p}\delta_n,\qquad P=\delta_0.
\]
对任意有界连续$\varphi$，远端项的贡献被$n^{-p}\|\varphi\|_\infty$控制并趋于零，故$P_n\Rightarrow P$；每一项也都有有限$p$阶矩。
但把远端质量搬回零点的成本为
\[
 W_p(P_n,P)^p=n^{-p}\cdot n^p=1,
\]
不趋于零。
这个反例把“远端质量很少”与“搬运成本很小”分开：缺少共同的距离尾部控制时，小概率乘大距离仍可保持常数。
涉及映射后的分布时同样要先检查输出矩，不能只因输入矩有限就使用目标空间的Wasserstein判据。

\subsection{弱收敛的判据与稳定性}

弱收敛定义只直接控制有界连续测试。
本节分别问事件概率怎样读取、连续变换后的测度是否仍收敛、无界成本能保留什么不等式。
这些是同一极限方式下的可用操作，不是与弱收敛并列的新收敛类型。

\subsubsection{事件概率与边界质量}

在Polish空间的Borel概率范围内，\term{Portmanteau定理}给出常用等价形式：
\[
 P_n\Rightarrow P
\]
等价于对所有闭集$C$有
\[
 \limsup_{n}P_n(C)\leq P(C),
\]
也等价于对所有开集$G$有
\[
 \liminf_nP_n(G)\geq P(G).
\]
由此，若Borel集$A$满足边界零质量$P(\partial A)=0$，则
\[
 P_n(A)\to P(A).
\]
这里只展开“从弱收敛推出这些可用界”的方向，整套等价形式作为标准定理使用。

对非空闭集$C$，取
\[
 h_k(x)=\max\{1-k\,d(x,C),0\}.
\]
每个$h_k$有界连续，且随$k$递减到$\mathbf 1_C$。
固定$k$时，由$\mathbf 1_C\leq h_k$和弱收敛，
\[
 \limsup_nP_n(C)\leq\lim_n\int h_k\,dP_n=\int h_k\,dP.
\]
再对固定极限测度用支配收敛令$k\to\infty$，得到闭集上界$\limsup_nP_n(C)\leq P(C)$；空闭集单独为零。
开集下界由补集得到：
\[
 \liminf_nP_n(G)\geq P(G).
\]
对一般集合$A$，用内部$A^\circ$和闭包$\overline A$夹住：
\[
 P(A^\circ)\leq\liminf_nP_n(A)\leq\limsup_nP_n(A)\leq P(\overline A).
\]
只有边界质量为零时两端才相等。
回看$\delta_{1/n}$与事件$\{0\}$：极限测度在边界$\{0\}$上有质量$1$，不满足条件，所以$P_n(\{0\})$不收敛到$P(\{0\})$。
事件收敛的例外出现在测试的不连续位置，不能只用“弱收敛比较宽松”一句话带过。

\subsubsection{连续推前与边缘保持}

设$T:X\to Y$是固定连续映射，$P_n\Rightarrow P$。
对任意目标端有界连续函数$\varphi$，先用推前积分公式改写：
\[
 \int_Y\varphi\,d(T_\#P_n)=\int_X\varphi\circ T\,dP_n.
\]
复合$\varphi\circ T$仍有界连续，由弱收敛右端趋于$\int_X\varphi\circ T\,dP$，所以
\[
 T_\#P_n\Rightarrow T_\#P.
\]
连续映射保持弱极限；这一步明确了在哪个空间用哪类测试，不能把结论推广到任意\term{可测映射}。

可测但不连续的映射可能失败。
在实轴上取$P_n=\delta_{1/n}$、$P=\delta_0$和$T(x)=\mathbf 1_{\{0\}}(x)$，则$T_\#P_n=\delta_0$而$T_\#P=\delta_1$，并不弱收敛。
连续坐标投影是乘积空间上的连续映射，所以弱收敛\term{耦合}列的边缘继承相应边缘极限；这是后面取得最优计划时保持边缘约束的依据。

\subsubsection{\term{下半连续}成本的极限下界}

对非负下半连续函数$c:X\to[0,\infty]$，可以取非负有界连续函数从下方逼近：
\[
 c_k(x)=\min\left\{k,\ \inf_y[c(y)+k\,d(x,y)]\right\}.
\]
每个$c_k$有界且连续（Lipschitz常数为$k$），随$k$递增并不超过$c$。
若在某点$x$有$\sup_kc_k(x)<c(x)$，取近似极小点$y_k$，则$d(x,y_k)\to0$而$y_k\to x$，下半连续性给出$c(x)\leq\liminf_kc(y_k)\leq\sup_kc_k(x)$，矛盾；故$c_k\uparrow c$。
若$c$处处为无穷，直接取$c_k=k$即可。

现在设$P_n\Rightarrow P$。先固定$k$，由$c_k\leq c$和弱收敛，
\[
 \int c_k\,dP=\lim_n\int c_k\,dP_n\leq\liminf_n\int c\,dP_n.
\]
再对极限测度使用单调收敛并令$k\to\infty$，得到下半连续积分界
\begin{equation}
 \int c\,dP\leq\liminf_n\int c\,dP_n.
 \label{eq:measure-lsc-portmanteau}
\end{equation}
这里测度随$n$变化，不能直接用固定测度的Fatou不等式替代整个论证；允许成本无界也不意味着积分取等号。

这条下界正是运输存在性论证中的成本步骤：接近最优的耦合列取弱子列后，极限成本不超过原列成本的下极限。
紧子列条件在下一小节核对；非空、有限和最小值达到仍是三个不同判断。

\subsection{\term{一致紧性}与弱子列}
\label{sec:measure-tightness}

前面先给定候选极限，再判断某种收敛是否成立。
本小节反过来问：一族概率是否至少留下一个可用的极限子列？
关键是能否用共同紧集把所有测度的大部分质量留在有界范围内；最后把这一判断用于\term{运输计划}的极小化序列。

\paragraph{一致紧性与极限子列}

对概率族$\mathcal F$定义一致紧性：对每个$\varepsilon>0$，都存在同一个紧集$K_\varepsilon$，使所有$P\in\mathcal F$满足
\[
 P(X\setminus K_\varepsilon)<\varepsilon.
\]
紧集可以随精度变化，但不能随族中的测度变化。
在无限维或一般度量空间中，有界球不必紧，统一矩界也不能无条件替代这一要求。

\begin{theorem}[\term{Prokhorov定理}，子列形式]
\label{thm:measure-prokhorov}
在Polish空间上，概率族$\mathcal F$一致紧，当且仅当$\mathcal F$中的每个序列都有弱收敛到某个Borel概率的子列。
\end{theorem}

极限概率不必属于原来的族$\mathcal F$；“定义某种收敛方式”与“保证存在极限子列”是两个判断。
这里说明使用条件，不展开完整证明。
实轴上的族$\{\delta_n:n\geq1\}$给出质量逃逸的典型反例：每个单项都有紧支撑$\{n\}$，但任何共同紧集都无法留住全部质量，序列没有弱收敛到概率的子列。
它与$P_n=(1-n^{-p})\delta_0+n^{-p}\delta_n$的矩尾反例分别说明“存在弱子列”与“几何成本趋零”需要不同控制，不能混为同一条件。

\paragraph{固定边缘下的最优计划}

在两个Polish空间上固定Borel概率$P,Q$。
先说明整个耦合族$\Pi(P,Q)$一致紧。
给定$\varepsilon>0$，由$P,Q$各自一致紧（单个概率在Polish空间上必一致紧），取紧集$K,L$使
\[
 P(K^c)<\frac\varepsilon2,\qquad Q(L^c)<\frac\varepsilon2.
\]
对任意$\gamma\in\Pi(P,Q)$，边缘约束给出
\[
 \gamma((K\times L)^c)
  \leq\gamma(K^c\times Y)+\gamma(X\times L^c)
  =P(K^c)+Q(L^c)<\varepsilon.
\]
所以是整族耦合共用同一个紧集$K\times L$，不是只证明某个候选单独紧。

设成本$c:X\times Y\to[0,\infty]$非负下半连续，并且至少有一个有限成本耦合。
令最优值
\[
 I=\inf_{\gamma\in\Pi(P,Q)}\int c\,d\gamma<\infty.
\]
按下确界选取耦合列$\gamma_n$使$\int c\,d\gamma_n\to I$。
由耦合族一致紧和Prokhorov定理，取弱收敛子列$\gamma_{n_k}\Rightarrow\gamma$。
连续坐标投影保持弱极限，所以$\gamma$的两个边缘仍为$P,Q$，即$\gamma\in\Pi(P,Q)$。
再由下半连续积分界~\eqref{eq:measure-lsc-portmanteau}，
\[
 \int c\,d\gamma\leq\liminf_k\int c\,d\gamma_{n_k}=I.
\]
而$\gamma$可行又给出$\int c\,d\gamma\geq I$，故最小值在$\gamma$处达到。
这就把紧性、边缘保持和成本下界三步接齐。
结论不保证最优确定性映射存在；下半连续成本的论证也不能用于任意可测成本。

\subsection{测度曲线与参数极限}

\paragraph{由序列判据读取时间连续性}

对固定空间上的概率曲线$t\mapsto\mu_t$，称它弱连续，若每个有界连续测试的读数$t\mapsto\int\varphi\,d\mu_t$连续；在时间区间上等价于
\[
 t_n\to t\ \Longrightarrow\ \mu_{t_n}\Rightarrow\mu_t.
\]
讨论$W_p$连续性时限定各时刻$\mu_t\in\mathcal P_p(X)$，并要求
\[
 W_p(\mu_{t_n},\mu_t)\to0,
\]
由矩判据~\eqref{eq:measure-wasserstein-criterion}，这还同时控制距离尾部。

在实轴上取$\mu_t=\delta_t$，则
\[
 W_p(\mu_s,\mu_t)=|s-t|,
\]
曲线关于$W_p$连续；但$s\neq t$时$\operatorname{TV}(\mu_s,\mu_t)=1$，曲线在全变差意义下处处不连续。
曲线是否连续取决于检测方式。
若每个时刻的测度原本位于不同空间，必须先固定到共同空间的表示，不能在不同底集上直接沿用同一个测试函数或距离。

\paragraph{固定时刻、有限时域与长时间}

对目标曲线$\mu_t$和近似曲线$\mu_t^h$，要区分三种越来越强的趋同说法：
\begin{enumerate}
 \item 每个固定时刻$t$，误差$W_p(\mu_t^h,\mu_t)\to0$；
 \item 每个有限时间区间$[0,T]$上，$\sup_{0\leq t\leq T}W_p(\mu_t^h,\mu_t)\to0$；
 \item $t\to\infty$时，曲线本身是否存在时间极限。
\end{enumerate}
三种量词不能互相替代。

在实轴上取$h>0$、$t\geq0$，令
\[
 \mu_t=\delta_0,\qquad \mu_t^h=\delta_{ht}.
\]
固定$t$时误差为$ht\to0$；有限时域$[0,T]$上的上界为$hT\to0$；但在$t=1/h$这一随$h$移动的时刻误差恒为$1$，不存在关于$h$的全域一致收敛。
这个短例保留了双极限边界。
演化方程如何产生测度曲线、数值近似为何稳定以及何时能交换极限，都留给后续相应问题。


\FloatBarrier
\section{本章小结}

本章把实数和欧氏空间上的长度、体积与积分推广到一般底集，得到四条主要认识。

第一，一般底集上必须另外给定可测范围、赋权和（在需要几何结论时）拓扑结构。
$\sigma$代数规定哪些集合可以测量，生成族使可测性和测度唯一性只需在简单集合上检验；测度规定这些集合的非负权重，有限、概率和$\sigma$有限是不同强度的质量条件；Borel结构、Polish空间和支持把可测集合与开闭邻域联系起来。
分层积分以及单调收敛、Fatou和支配收敛的机制仍沿用微积分，乘积结构、Tonelli/Fubini和柱集接口则把积分推广到联合对象。

第二，固定同一空间时，零集关系说明两项赋权是否相容。
绝对连续、等价和相互奇异分别刻画零集的单向蕴含、双向共享和质量范围分离；Radon–Nikodym密度恢复绝对连续部分的权重，Lebesgue分解把不能用密度表示的奇异部分另外拆出；有符号测度的Hahn–Jordan分解和总变差累计正负变化，不允许它们抵消。
全变差、Hellinger、KL和Wasserstein从不同机制量化概率差异：前三者在固定位置比较事件或密度，Wasserstein还要使用共同距离；积分概率度量把这些指标统一成“某类测试函数的最大积分差”，而不是再添加一个互斥的指标类别。

第三，跨空间时，可测映射为每个输入指定一个目标点，测度核为每个输入指定一项目标测度。
在输入测度有限、核有统一有限质量界的范围内，核积分给出输出测度并保持观测传递，核复合扩展了映射复合；联合测度同时记录输入和输出，乘积和图测度都是其特例，沿第一边缘在标准Borel条件下可反向分解出归一化核。
核的总质量可以随输入变化，归一化条件才保证质量保持。
较粗可测结构上的条件积分建立了通用的积分保持关系；条件期望的代数法则、Bayes公式和随机过程解释由后文概率论继续展开。

第四，建立共同空间后比较的是输出权重，给定跨空间成本后优化的是配对代价。
同一归一化测度核收缩全变差，但这一结论不能无条件推广到两个不同的传递规则；表示后的零差异只说明该表示看不到某些输入区别，不能恢复丢掉的信息。
跨空间运输要求先固定可测的成本函数；Kantorovich计划允许拆分质量，Monge映射不允许，最优计划达到还要依赖一致紧性和成本下半连续。
比较之前应先固定对应关系、允许的观测和成本，不能只凭两侧内部距离就直接计算跨空间距离。

第五，序列与曲线的极限依赖所用标准。
逐集、全变差、弱收敛和Wasserstein收敛依次放宽或加入不同结构；弱收敛由有界连续测试定义，Portmanteau界控制边界零质量事件，连续推前保持弱极限，下半连续成本只给积分下界，一致紧性经Prokhorov定理保证弱子列，Wasserstein收敛还要额外控制距离尾部。
后文的随机过程、条件模型、经验测度和策略分析都可以调用这些通用操作，但它们的概率性质、统计保证和生成机制仍需各自的条件。


